\documentclass[
	% -- opções da classe memoir --
	12pt,				% tamanho da fonte
	openright,			% capítulos começam em pág ímpar (insere página vazia caso preciso)
	twoside,			% para impressão em recto e verso. Oposto a oneside
	a4paper,			% tamanho do papel.
	% -- opções do pacote babel --
	english,			% idioma adicional para hifenização
	brazil				% o último idioma é o principal do documento
	]{abntex2}

% ---
% Pacotes básicos do modelo abnTeX2
% ---
\usepackage{lmodern}			% Usa a fonte Latin Modern
\usepackage[T1]{fontenc}		% Seleção de códigos de fonte.
\usepackage[utf8]{inputenc}
% Codificação do documento (conversão automática dos acentos)
\usepackage{indentfirst}		% Indenta o primeiro parágrafo de cada seção.
\usepackage{color}				% Controle das cores
\usepackage[table,xcdraw]{xcolor}
\usepackage{graphicx}			% Inclusão de gráficos
\usepackage{microtype} 			% Para melhorias de justificação
\usepackage{amsmath}			% Ambientes matemáticos formais
\usepackage{booktabs}			% Tabelas no padrão IBGE (sem linhas verticais)
\usepackage{multirow}			% Mesclagem de linhas em tabelas/quadros
\usepackage{array}				% Formatação avançada de colunas em tabelas/quadros
\usepackage[acronym,nonumberlist]{glossaries}
% Suporte aos comandos \acrshort, \acrlong, \gls

% ---
% Pacotes de citações padrão abnTeX2
% ---
\usepackage[brazilian,hyperpageref]{backref}	% Páginas com as citações na bibliografia
\usepackage[alf]{abntex2cite}					% Citações padrão ABNT (autor-data)

% ---
% Configurações do pacote backref (padrão abnTeX2)
% ---
\renewcommand{\backrefpagesname}{Citado na(s) página(s):~}
\renewcommand{\backref}{}
\renewcommand*{\backrefalt}[4]{%
	\ifcase #1 %
		Nenhuma citação no texto.%
	\or
		Citado na página #2.%
	\else
		Citado #1 vezes nas páginas #2.%
	\fi}%

% ---
% Caminhos de imagens e contador auxiliar de Quadros
% ---
\graphicspath{{./}{docs/latex/TCC_II/}}
\newcounter{quadro}

% ---
% Espaçamentos entre linhas e parágrafos (padrão abnTeX2)
% ---
\setlength{\parindent}{1.3cm}
\setlength{\parskip}{0.2cm}
\emergencystretch=1.5em

% ---
% Definições de siglas
% ---
\makeglossaries
\glsdisablehyper
\newacronym{acid}{ACID}{Atomicity, Consistency, Isolation, Durability}
\newacronym{acl}{ACL}{Access Control List}
\newacronym{api}{API}{Application Programming Interface}
\newacronym{asgi}{ASGI}{Asynchronous Server Gateway Interface}
\newacronym{can}{CAN}{Controller Area Network}
\newacronym{cba}{CBA}{Component-Based Architecture}
\newacronym{ci}{CI}{Circuito Integrado}
\newacronym{crispdm}{CRISP-DM}{Cross-Industry Standard Process for Data Mining}
\newacronym{dlc}{DLC}{Data Length Code}
\newacronym{dtc}{DTC}{Diagnostic Trouble Codes}
\newacronym{ecu}{ECU}{Electronic Control Unit}
\newacronym{espidf}{ESP-IDF}{Espressif IoT Development Framework}
\newacronym{gpio}{GPIO}{General Purpose Input/Output}
\newacronym{gpos}{GPOS}{General-Purpose Operating System}
\newacronym{hal}{HAL}{Hardware Abstraction Layer}
\newacronym{http}{HTTP}{Hypertext Transfer Protocol}
\newacronym{ia}{IA}{Inteligência Artificial}
\newacronym{iag}{IAG}{Inteligência Artificial Generativa}
\newacronym{ihc}{IHC}{Interface Humano-Computador}
\newacronym{imu}{IMU}{Unidade de Medição Inercial}
\newacronym{iot}{IoT}{Internet of Things}
\newacronym{iso}{ISO}{International Organization for Standardization}
\newacronym{isr}{ISR}{Interrupt Service Routine}
\newacronym{json}{JSON}{JavaScript Object Notation}
\newacronym{mitm}{MITM}{Man-in-the-Middle}
\newacronym{ml}{ML}{Machine Learning}
\newacronym{mqtt}{MQTT}{Message Queuing Telemetry Transport}
\newacronym{mqtts}{MQTTS}{Message Queuing Telemetry Transport Secure}
\newacronym{obd}{OBD-II}{On-Board Diagnostics Version II}
\newacronym{orm}{ORM}{Object-Relational Mapping}
\newacronym{osi}{OSI}{Open Systems Interconnection}
\newacronym{pc}{PC}{Personal Computer}
\newacronym{pid}{PID}{Parameters Identification}
\newacronym{rpm}{RPM}{Rotações por Minuto}
\newacronym{rtos}{RTOS}{Real-Time Operating System}
\newacronym{sae}{SAE}{Society of Automotive Engineers}
\newacronym{sd}{SD}{Secure Digital}
\newacronym{smp}{SMP}{Symmetric Multiprocessing}
\newacronym{sql}{SQL}{Structured Query Language}
\newacronym{ssot}{SSOT}{Single Source of Truth}
\newacronym{tcc}{TCC}{Trabalho de Conclusão de Curso}
\newacronym{tls}{TLS}{Transport Layer Security}
\newacronym{twai}{TWAI}{Two-Wire Automotive Interface}
\newacronym{uart}{UART}{Universal Asynchronous Receiver-Transmitter}
\newacronym{uce}{UCE}{Unidade de Controle Eletrônico}
\newacronym{wifi}{Wi-Fi}{Wireless Fidelity}


% ---
% Informações de dados para CAPA e FOLHA DE ROSTO (abnTeX2)
% ---
\titulo{Telemetria Automotiva Inteligente: Sistema Embarcado via OBD-II/CAN Integrado à
Nuvem para Detecção Não Supervisionada de Anomalias}
\autor{Lucas Praxedes Fischer de Mattos}
\local{São José dos Campos}
\data{2026}
\orientador{Prof. Dr. André Marcorin de Oliveira}
\instituicao{%
  Universidade Federal de São Paulo -- UNIFESP
  \par
  Instituto de Ciência e Tecnologia
  \par
  Bacharelado em Engenharia da Computação}
\tipotrabalho{Trabalho de Conclusão de Curso (Graduação)}
\preambulo{Trabalho de Conclusão de Curso apresentado ao Instituto de Ciência e
Tecnologia da Universidade Federal de São Paulo -- Campus São José dos Campos, como
exigência parcial para obtenção do título de Bacharel em Engenharia da Computação.}

% ---
% Configurações de aparência do PDF final
% ---
\makeatletter
\hypersetup{
     	pdftitle={\@title},
	    pdfauthor={\@author},
    	pdfsubject={\imprimirpreambulo},
	    pdfcreator={LaTeX with abnTeX2},
		colorlinks=true,
    	linkcolor=black,
    	citecolor=black,
    	filecolor=black,
		urlcolor=black,
		bookmarksdepth=4
}
\makeatother

% ---
% Início do documento
% ---
\begin{document}

\frenchspacing

% =========================================================================
% ELEMENTOS PRÉ-TEXTUAIS
% =========================================================================
\imprimircapa

\imprimirfolhaderosto*

% ---
% Declaração sobre o uso de Inteligência Artificial Generativa (IN nº 06/2026 UNIFESP)
% ---
\pretextualchapter{Declaração sobre o uso de Inteligência Artificial Generativa}

Durante a elaboração deste trabalho, o autor utilizou a ferramenta de Inteligência
Artificial Generativa (\acrshort{iag}) Google Gemini (integrada ao ambiente de
desenvolvimento) com a finalidade de apoio na revisão linguística e gramatical,
organização estrutural de seções previamente concebidas pelo autor e verificação
exploratória de lacunas técnicas (\textit{gap analysis}) e casos de borda na
documentação arquitetural e na suíte de testes do projeto. Após esse uso, o autor
revisou, validou empiricamente e editou o conteúdo conforme necessário, assumindo
integral e exclusiva responsabilidade intelectual e técnica pelo conteúdo do trabalho.
\cleardoublepage

% ---
% RESUMO e ABSTRACT (ABNT NBR 6028 / Alto Nível Arquitetural)
% ---
\setlength{\absparsep}{18pt}
\begin{resumo}
A crescente complexidade dos sistemas eletrônicos automotivos evidencia a importância da
coleta e análise contínua de dados operacionais dos veículos; contudo, o acesso a essas
informações em tempo real permanece restrito em automóveis comerciais de entrada, ao
passo que ferramentas convencionais de oficina limitam-se à leitura reativa de falhas já
consolidadas. Este trabalho apresenta o desenvolvimento de um sistema computacional
embarcado integrado à nuvem via barramento \textit{Controller Area Network} (CAN) e
padrão \textit{On-Board Diagnostics Version II} (OBD-II), voltado ao monitoramento
telemétrico ativo e à detecção não supervisionada de anomalias operacionais. A
metodologia divide-se na construção de uma unidade de aquisição embarcada baseada em
microcontrolador de núcleo duplo operando sobre sistema operacional de tempo real sob
arquitetura modular em linguagem C para leitura periódica de rotação do motor,
velocidade veicular e temperatura do líquido de arrefecimento, aliada à transmissão sem
fio via protocolo de mensageria seguro com criptografia na camada de transporte. Na
infraestrutura de servidor, implementou-se um serviço de aplicação assíncrono
responsável pela validação de contratos de dados, persistência transacional em banco de
dados relacional e disponibilização das séries temporais em um painel de visualização
interativo em tempo real, complementado pela modelagem analítica de um fluxo de
aprendizado de máquina não supervisionado fundamentado em árvores de isolamento para
identificação proativa de desvios térmicos e cinemáticos no funcionamento do motor.

\textbf{Palavras-chave}: telemetria veicular. OBD-II. sistemas embarcados. internet das
coisas. detecção de anomalias.
\end{resumo}

\begin{resumo}[Abstract]
\begin{otherlanguage*}{english}
The growing complexity of automotive electronic systems highlights the importance of
continuous collection and analysis of vehicle operational data; however, real-time
access to this information remains restricted in entry-level commercial vehicles, while
conventional diagnostic tools are largely limited to reactive fault reading. This work
presents the development of a cloud-integrated embedded computing system using the
Controller Area Network (CAN) bus and the On-Board Diagnostics Version II (OBD-II)
standard, aimed at active telemetric monitoring and unsupervised detection of
operational anomalies. The methodology is structured around an embedded acquisition unit
based on a dual-core microcontroller running a real-time operating system under a
modular architecture in C language for periodic reading of engine speed, vehicle speed,
and coolant temperature, coupled with wireless transmission via a secure messaging
protocol with transport-layer encryption. On the server infrastructure, an asynchronous
application service was implemented to perform data contract validation, transactional
persistence in a relational database, and real-time visualization through an interactive
dashboard, complemented by the analytical modeling of an unsupervised machine learning
pipeline based on isolation trees for proactive identification of thermal and kinematic
deviations in engine operation.

\vspace{\onelineskip}

\noindent
\textbf{Keywords}: vehicular telemetry. OBD-II. embedded systems. internet of things.
anomaly detection.
\end{otherlanguage*}
\end{resumo}

% ---
% Lista de ilustrações
% ---
\pdfbookmark[0]{\listfigurename}{lof}
\listoffigures*
\cleardoublepage

% ---
% Lista de tabelas
% ---
\pdfbookmark[0]{\listtablename}{lot}
\listoftables*
\cleardoublepage

% ---
% Lista de abreviaturas e siglas (padrão abnTeX2)
% ---
\begin{siglas}
  \item[ACID] \textit{Atomicity, Consistency, Isolation, Durability}
  \item[ACL] \textit{Access Control List}
  \item[API] \textit{Application Programming Interface}
  \item[ASGI] \textit{Asynchronous Server Gateway Interface}
  \item[CAN] \textit{Controller Area Network}
  \item[CBA] \textit{Component-Based Architecture}
  \item[CI] Circuito Integrado
  \item[CRISP-DM] \textit{Cross-Industry Standard Process for Data Mining}
  \item[DLC] \textit{Data Length Code}
  \item[DTC] \textit{Diagnostic Trouble Codes}
  \item[ECU] \textit{Electronic Control Unit}
  \item[ESP-IDF] \textit{Espressif IoT Development Framework}
  \item[GPIO] \textit{General Purpose Input/Output}
  \item[GPOS] \textit{General-Purpose Operating System}
  \item[HAL] \textit{Hardware Abstraction Layer}
  \item[HTTP] \textit{Hypertext Transfer Protocol}
  \item[IA] Inteligência Artificial
  \item[IAG] Inteligência Artificial Generativa
  \item[IHC] Interface Humano-Computador
  \item[IMU] Unidade de Medição Inercial
  \item[IoT] \textit{Internet of Things}
  \item[ISO] \textit{International Organization for Standardization}
  \item[ISR] \textit{Interrupt Service Routine}
  \item[JSON] \textit{JavaScript Object Notation}
  \item[MITM] \textit{Man-in-the-Middle}
  \item[ML] \textit{Machine Learning}
  \item[MQTT] \textit{Message Queuing Telemetry Transport}
  \item[MQTTS] \textit{Message Queuing Telemetry Transport Secure}
  \item[OBD-II] \textit{On-Board Diagnostics Version II}
  \item[ORM] \textit{Object-Relational Mapping}
  \item[OSI] \textit{Open Systems Interconnection}
  \item[PC] \textit{Personal Computer}
  \item[PID] \textit{Parameters Identification}
  \item[RPM] Rotações por Minuto
  \item[RTOS] \textit{Real-Time Operating System}
  \item[SAE] \textit{Society of Automotive Engineers}
  \item[SD] \textit{Secure Digital}
  \item[SMP] \textit{Symmetric Multiprocessing}
  \item[SQL] \textit{Structured Query Language}
  \item[SSOT] \textit{Single Source of Truth}
  \item[TCC] Trabalho de Conclusão de Curso
  \item[TLS] \textit{Transport Layer Security}
  \item[TWAI] \textit{Two-Wire Automotive Interface}
  \item[UART] \textit{Universal Asynchronous Receiver-Transmitter}
  \item[UCE] Unidade de Controle Eletrônico
  \item[Wi-Fi] \textit{Wireless Fidelity}
\end{siglas}

% ---
% Sumário
% ---
\pdfbookmark[0]{\contentsname}{toc}
\tableofcontents*
\cleardoublepage

% =========================================================================
% ELEMENTOS TEXTUAIS
% =========================================================================
\textual
\glsresetall

\chapter{Introdução}

A indústria automotiva tem passado por uma profunda transformação impulsionada pela
demanda crescente por veículos mais inteligentes, conectados e eficientes
\cite{llopis2021}. Nesse cenário, a capacidade de coletar, processar e transmitir dados
veiculares em tempo real consolidou-se como um pilar tecnológico para a engenharia
automotiva moderna. A ampla gama de informações geradas pelos subsistemas eletrônicos
embarcados viabiliza desde o monitoramento contínuo da integridade termomecânica do
motor até o desenvolvimento de arquiteturas de manutenção preditiva e análise de perfil
de condução \cite{andria2016}.

No núcleo da comunicação interna dos automóveis contemporâneos encontra-se a rede
\acrshort{can} (do inglês, \textit{Controller Area Network}), um protocolo de barramento
serial diferencial que se destaca por sua robustez eletromagnética, determinismo na
arbitragem de mensagens e elevada confiabilidade operacional \cite{bosch2012}. Embora
essa infraestrutura de comunicação seja onipresente nas Unidades de Controle Eletrônico
(\acrshort{ecu}s), o acesso estruturado às grandezas físicas que trafegam no barramento
permanece restrito para o usuário final, sobretudo em veículos compactos de passeio,
cujos painéis de instrumentos suprimem indicadores detalhados de saúde térmica e
dinâmica do motor \cite{guimaraes2007}.

A modificação invasiva dos chicotes ou dos sistemas eletrônicos originais para extração
dessas grandezas mostra-se técnica e economicamente inviável, implicando riscos à
segurança operacional e perda de garantia de fábrica. Uma abordagem de engenharia eficaz
e não intrusiva consiste na utilização de analisadores de barramento (\textit{sniffers})
e dispositivos de aquisição telemétrica conectados à interface padronizada de
diagnóstico de bordo \acrshort{obd} (do inglês,
\textit{On-Board Diagnostics Version II}), capturando e interrogando os parâmetros
operacionais da \acrshort{ecu} no barramento \acrshort{can} sem interferir nas malhas de
controle críticas do veículo \cite{iso15765,csselectronics2024}.

A convergência entre esses sistemas embarcados de aquisição veicular, tecnologias da
Internet das Coisas (do inglês, \textit{Internet of Things} -- \acrshort{iot}) para
conectividade em nuvem e técnicas de aprendizado de máquina (\acrshort{ml}, do inglês
\textit{Machine Learning}) representa um avanço substancial frente aos leitores de
diagnóstico convencionais \cite{iqbal2019,singh2021}. Ao transmitir a
telemetria de borda sob canais criptografados para uma infraestrutura de persistência e
processamento em nuvem, torna-se possível superar o paradigma puramente reativo dos
códigos de falha tradicionais, viabilizando a modelagem de mecanismos de detecção não
supervisionada de anomalias operacionais antes que evoluam para panes mecânicas severas.

\section{Objetivo geral}

O presente trabalho tem como objetivo geral o desenvolvimento de um sistema
computacional embarcado integrado à nuvem, de baixo custo e alta confiabilidade, voltado
à telemetria automotiva ativa e à análise preditiva. O sistema é projetado para realizar
a leitura em tempo real dos parâmetros padronizados da norma \acrshort{sae} J1979 no
barramento \acrshort{can} (especificamente rotação do motor em \acrshort{rpm},
velocidade veicular e temperatura do líquido de arrefecimento), efetuar a transmissão
segura dos pacotes via protocolo \acrshort{mqtts} para um servidor de aplicação com
persistência relacional transacional, e disponibilizar o monitoramento em um
\textit{dashboard} interativo aliado à modelagem de detecção não supervisionada de
anomalias operacionais por meio de aprendizado de máquina (\acrshort{ml}).

\section{Objetivos específicos}

Para alcançar o objetivo geral proposto, estabeleceram-se os seguintes objetivos
específicos:

\begin{itemize}
    \item Projetar e integrar o circuito de aquisição embarcado baseado no
        microcontrolador ESP32 (utilizando seu controlador nativo \acrshort{twai}) e no
        transceptor \acrshort{can} SN65HVD230 operando a $500\text{ kbaud}$.
    \item Desenvolver o \textit{firmware} embarcado em linguagem C nativa sobre o
        \textit{framework} ESP-IDF e sistema operacional de tempo real (\acrshort{rtos})
        FreeRTOS, estruturado sob Arquitetura Baseada em Componentes (\acrshort{cba}).
    \item Implementar a camada de transporte seguro via \acrshort{mqtts} sobre
        \acrshort{tls} v1.2 e o \textit{backend} de ingestão assíncrona em Python
        (utilizando FastAPI, Pydantic e SQLAlchemy/SQLite).
    \item Construir uma aplicação \textit{web} interativa em Streamlit para
        monitoramento telemétrico em tempo real e validação por emulação sintética.
    \item Estruturar e modelar a arquitetura do \textit{pipeline} de aprendizado de
        máquina fundamentado na metodologia \acrshort{crispdm} e no algoritmo
        \textit{Isolation Forest} para detecção não supervisionada de anomalias no
        funcionamento do motor.
\end{itemize}

\chapter{Fundamentação Teórica}

Para o desenvolvimento deste trabalho, foram investigados os fundamentos de sistemas
embarcados automotivos, redes de comunicação veicular sob o padrão \acrshort{obd},
sistemas operacionais de tempo real, protocolos seguros de mensageria para
\acrshort{iot} e métodos de aprendizado de máquina para detecção de anomalias, os quais
são formalizados nas seções a seguir.

\section{Unidade de Controle Eletrônico}

Uma Unidade de Controle Eletrônico (\acrshort{uce}) automotiva, consagrada na literatura
internacional pela sigla \acrshort{ecu} (do inglês, \textit{Electronic Control Unit}),
consiste em um sistema computacional embarcado dotado de microcontrolador e memórias não
voláteis onde residem as rotinas responsáveis pelo gerenciamento de subsistemas
específicos do veículo \cite{guimaraes2007,andrijauskas2018}.

Nesse sentido, uma \acrshort{ecu} opera de maneira análoga a um computador dedicado de
tempo real projetado sob restrições estritas de custo, consumo energético e tempo de
resposta \cite{peckol2008,ganssle2008}. Internamente, sua placa de circuito impresso
integra unidades de processamento, conversores analógico-digitais e interfaces de
potência, executando o \textit{software} embarcado ou \textit{firmware} de controle e
autodiagnóstico dos subsistemas mecânicos e elétricos aos quais se encontra acoplada.

Com a evolução das exigências de eficiência energética, segurança ativa e controle de
emissões, as montadoras intensificaram a distribuição de módulos eletrônicos pelo
veículo, abrangendo o controle da injeção eletrônica, transmissão, freios antibloqueio,
climatização e carroceria, conforme ilustrado na Figura~\ref{fig:uces_automovel}. A
\acrshort{ecu} principal do motor coordena malhas fechadas de injeção de combustível,
avanço de ignição, controle de mistura estequiométrica e monitoramento de falhas por
meio de múltiplos sensores e atuadores distribuídos pelo conjunto motriz
\cite{capelli2010,guimaraes2007}.

\begin{figure}[!htb]
	\centering
	\caption{Unidades de Controle Eletrônico (UCEs) de um automóvel}
	\includegraphics[width=0.78\textwidth]{imagens/uces_automovel_capelli.jpeg}
	\legend{Fonte: \citeonline{capelli2010} \textit{apud} \citeonline{yu2024}.}
	\label{fig:uces_automovel}
\end{figure}

\section{A rede de comunicação veicular: CAN Bus}

O protocolo \acrshort{can} foi desenvolvido na década de 1980 pela Robert Bosch GmbH e
consolidou-se como padrão internacional para a comunicação serial entre as diversas
\acrshort{ecu}s presentes em um veículo \cite{bosch2012}. Sua adoção em larga escala
pela indústria automotiva decorre de sua elevada confiabilidade, arbitragem não
destrutiva de barramento e imunidade a interferências eletromagnéticas, requisitos
críticos no ambiente veicular. A topologia do \acrshort{can} \textit{Bus} baseia-se em
um barramento compartilhado de par trançado ao qual múltiplos nós se conectam
paralelamente, reduzindo de modo expressivo a complexidade e a massa dos chicotes
elétricos em comparação às ligações ponto a ponto \cite{guimaraes2007}.

Na camada física, a sinalização da rede \acrshort{can} de alta velocidade opera por
transmissão diferencial através de duas linhas: \acrshort{can} \textit{High}
(\texttt{CAN\_H}) e \acrshort{can} \textit{Low} (\texttt{CAN\_L}), conforme representado
na Figura~\ref{fig:sinais_can}. Os dados lógicos são codificados em dois estados
elétricos: dominante (\textit{bit} \texttt{0}) e recessivo (\textit{bit} \texttt{1}). No
estado recessivo, o potencial elétrico em ambas as linhas repousa em aproximadamente
$2{,}5\text{ V}$; no estado dominante, a tensão em \texttt{CAN\_H} eleva-se para cerca
de $3{,}5\text{ V}$ enquanto em \texttt{CAN\_L} reduz-se para aproximadamente
$1{,}5\text{ V}$, resultando em uma diferença de potencial nominal de $2{,}0\text{ V}$
\cite{iso11898,csselectronics2023}. Essa característica diferencial assegura que ruídos
eletromagnéticos induzidos de modo comum afetem ambas as linhas simultaneamente,
preservando a diferença de tensão e a integridade da informação.

\begin{figure}[!htb]
	\centering
	\caption{Sinais elétricos da transmissão CAN}
	\includegraphics[width=0.75\textwidth]{imagens/sinais_eletricos_can_vernizzi.jpeg}
	\legend{Fonte: Adaptado de \citeonline{embarcados2020} \textit{apud}
	\citeonline{vernizzi2024}.}
	\label{fig:sinais_can}
\end{figure}

O empacotamento dos \textit{bits} transmitidos compõe quadros estruturados denominados
\acrshort{can} \textit{frames}, os quais carregam o campo de início de quadro, o
identificador de arbitragem (de $11\text{ bits}$ no formato padrão ou $29\text{ bits}$
no formato estendido), o código de comprimento de dados (\acrshort{dlc}, do inglês
\textit{Data Length Code}), o campo de carga útil (\textit{payload}) de até
$8\text{ bytes}$ e o verificador cíclico de redundância trocado entre as \acrshort{ecu}s
\cite{bosch2012,iso11898}.

\section{Diagnóstico de bordo: o padrão OBD-II}

O sistema de Diagnóstico de Bordo (do inglês, \textit{On-Board Diagnostics}) surgiu da
necessidade regulatória de monitorar continuamente os subsistemas veiculares associados
ao controle de emissões e à integridade do motor. A sua segunda geração, denominada
\acrshort{obd}, tornou-se obrigatória nos Estados Unidos a partir de 1996 e,
posteriormente, foi padronizada internacionalmente, inclusive no mercado brasileiro
\cite{guimaraes2007}.

O padrão \acrshort{obd} especifica um conector físico trapezoidal de 16 pinos definido
pela norma \acrshort{sae} J1962, posicionado no habitáculo próximo ao condutor como
interface universal de diagnóstico, cuja disposição de pinos e protocolos associados é
apresentada na Figura~\ref{fig:obd2_plug}. Por meio desse conector, torna-se possível
requisitar grandezas operacionais em tempo real, tais como velocidade veicular, rotação
do motor (\acrshort{rpm}), temperatura do líquido de arrefecimento e Códigos de Falha de
Diagnóstico (\acrshort{dtc}, do inglês \textit{Diagnostic Trouble Codes})
\cite{guimaraes2007,saej1979}.

\begin{figure}[!htb]
	\centering
	\caption{Pinagem do conector padronizado OBD-II (SAE J1962) e protocolos de
	comunicação}
	\includegraphics[width=0.72\textwidth]{imagens/obd2 plug.png}
	\legend{Fonte: Adaptado de \citeonline{guimaraes2007}.}
	\label{fig:obd2_plug}
\end{figure}

Segundo \citeonline{guimaraes2007}, o autodiagnóstico fundamenta-se no registro de
códigos \acrshort{dtc} na memória dos módulos eletrônicos sempre que leituras de
sensores extrapolam os limites operacionais projetados pelo fabricante. Na camada de
comunicação, a especificação \acrshort{obd} admite cinco protocolos distintos conforme
os pinos povoados no conector (Figura~\ref{fig:obd2_plug}): \acrshort{sae} J1850 PWM
(pinos 2 e 10), \acrshort{sae} J1850 VPW (pino 2), ISO 9141-2 / ISO 14230-4
\textit{K-Line} (pinos 7 e 15) e o padrão ISO 15765-4 baseado em barramento diferencial
\acrshort{can} nos pinos 6 (\texttt{CAN\_H}) e 14 (\texttt{CAN\_L}), o qual se
consolidou como o protocolo de alta velocidade obrigatório na frota fabricada a partir
de 2008 \cite{iso15765,guimaraes2007}.

\subsection{Convergência entre o padrão OBD-II e o protocolo CAN no modelo
OSI}\label{sec:convergencia_osi}

Historicamente, a convergência entre o padrão \acrshort{obd} e o protocolo
\acrshort{can} decorreu da necessidade de superar a fragmentação tecnológica da década
de
1990, quando cada grupo automotivo adotava barramentos de diagnóstico distintos e
incompatíveis entre si (como o \acrshort{sae} J1850 PWM na Ford, o \acrshort{sae} J1850
VPW na General Motors e as variações ISO 9141-2 / ISO 14230-4 \textit{K-Line} em
montadoras europeias e asiáticas), operando em baixas taxas de transmissão de
$10{,}4\text{ kbaud}$ a $41{,}6\text{ kbaud}$ \cite{guimaraes2007}. Sob a ótica prática
de engenharia, à medida que o barramento \acrshort{can} já se consolidava internamente
nos veículos para interligar as \acrshort{ecu}s de trem de força e segurança ativa em
tempo real, manter uma segunda rede física exclusiva e lenta apenas para diagnóstico
externo tornou-se técnica e economicamente redundante. A unificação regulatória via
norma ISO 15765-4 (\textit{Diagnostics on Controller Area Networks} --- DoCAN) permitiu
reaproveitar a própria infraestrutura diferencial de par trançado do veículo (pinos 6 e
14) operando a $250\text{ kbaud}$ ou $500\text{ kbaud}$, elevando a largura de banda de
diagnóstico em até $48\text{ vezes}$ frente à \textit{K-Line} e garantindo arbitragem
determinística sem perda de pacotes \cite{iso15765}.

Para compreender de forma rigorosa como o \acrshort{can} e o \acrshort{obd} operam em
conjunto sem sobreposição de funções, recorre-se ao modelo de referência \acrshort{osi}
(\textit{Open Systems Interconnection}), estruturado em sete camadas hierárquicas de
comunicação nas quais cada nível presta serviços à camada imediatamente superior por
meio do encapsulamento sucessivo de unidades de dados, conforme ilustrado na
Figura~\ref{fig:modelo_osi_can} \cite{tanenbaum2011,lugli2009}. A especificação original
do protocolo \acrshort{can} (ISO 11898) define exclusivamente as duas camadas inferiores
do modelo \acrshort{osi}: a Camada 1 (Física, ou \textit{Physical Layer}), que
estabelece os níveis elétricos diferenciais e a temporização de \textit{bits}, e a
Camada 2 (Enlace de Dados, ou \textit{Data Link Layer}), responsável pelo empacotamento
em quadros de até $8\text{ bytes}$, endereçamento por identificadores, controle de
acesso ao meio e verificação cíclica de erros, sem atribuir significado semântico ao
conteúdo da carga útil (\textit{payload}).

Enquanto as camadas intermediárias (Camadas 3 a 6) não são especificadas nativamente
pelo protocolo \acrshort{can} base e a Camada 7 (Aplicação, ou
\textit{Application Layer}) permanece aberta para definição pela aplicação
(Figura~\ref{fig:modelo_osi_can}), a arquitetura de diagnóstico automotivo preenche
exatamente essas lacunas. O padrão \acrshort{obd}, regido pela norma \acrshort{sae}
J1979, ocupa o topo da pilha na Camada 7 (Aplicação), padronizando o dicionário de
serviços de diagnóstico, os códigos de falha (\acrshort{dtc}s), os identificadores de
parâmetros (\acrshort{pid}s) e as equações de conversão em grandezas físicas, ao passo
que a norma ISO 15765 atua nas Camadas 3 e 4 (Rede e Transporte) como o elo de ligação
que encapsula e segmenta as mensagens de diagnóstico \acrshort{obd} no interior dos
quadros de enlace da rede \acrshort{can} \cite{iso15765,saej1979}.

\begin{figure}[!htb]
	\centering
	\caption{O protocolo CAN dentro do modelo OSI}
	\includegraphics[width=0.68\textwidth]{imagens/modelo_osi_can_farinelli.png}
	\legend{Fonte: Adaptado de \citeonline{lugli2009} \textit{apud}
	\citeonline{farinelli2016}.}
	\label{fig:modelo_osi_can}
\end{figure}

\section{Aquisição de dados via OBD-II}\label{sec:aquisicao_obd2}

A interação ativa com a \acrshort{ecu} para extração de telemetria em tempo real
baseia-se no envio de Identificadores de Parâmetros (\acrshort{pid}s, do inglês
\textit{Parameters Identification}), cujos modos de operação, códigos de serviço e
equações de escalonamento são padronizados pela norma \acrshort{sae} J1979
\cite{saej1979}. No Modo de Serviço \texttt{0x01} (exibição de dados
atuais do trem de força), cada \acrshort{pid} hexadecimal referencia uma variável física
específica, a exemplo da temperatura do líquido de arrefecimento (\acrshort{pid}
\texttt{0x05}), rotação do motor (\acrshort{pid} \texttt{0x0C}) e velocidade do veículo
(\acrshort{pid} \texttt{0x0D}) \cite{csselectronics2023}.

Em redes regidas pela norma ISO 15765-4 (barramento \acrshort{can} de $11\text{ bits}$ a
$500\text{ kbaud}$), o dispositivo leitor transmite um quadro de requisição utilizando o
identificador funcional de \textit{broadcast} \texttt{0x7DF} \cite{iso15765}. A
\acrshort{ecu} principal do motor processa a solicitação e retorna um quadro de resposta
no identificador físico \texttt{0x7E8}, contendo os bytes brutos de medição
($A, B, C, D$), os quais são submetidos às expressões algébricas da norma \acrshort{sae}
J1979 para conversão nas respectivas grandezas físicas de engenharia
\cite{sitohang2014}.

\section{Microcontroladores para aplicações IoT}\label{sec:microcontroladores}

A seleção da plataforma de processamento embarcado constitui uma decisão crítica no
projeto de unidades telemétricas, exigindo equilíbrio entre determinismo de leitura do
barramento automotivo, dimensionamento de memória, suporte nativo a pilhas
criptográficas de rede e custo reduzido \cite{peckol2008,andrijauskas2018}.
Diversas arquiteturas são reportadas na literatura:

\begin{itemize}
    \item \textbf{ATmega328P e Arduino UNO:} Empregados em protótipos elementares de
        aquisição de sensores e \textit{data logging}, porém limitados pela ausência de
        conectividade sem fio nativa, reduzida memória SRAM ($2\text{ KB}$) e ausência
        de suporte a sistemas operacionais multitarefa para pilhas de segurança modernas
        \cite{atmel_atmega328p,andrijauskas2018}.
    \item \textbf{Microcontroladores PIC e ARM Cortex-M:} Utilizados historicamente em
        interfaces dedicadas de diagnóstico automotivo (como os interpretadores baseados
        na família PIC18 de 8 bits) e em nós de rede \acrshort{can} de 32 bits
        \cite{microchip_pic18f4580,st_stm32f103}, mas frequentemente dependentes de
        módulos transceptores de rádio externos para integração direta com arquiteturas
        de nuvem.
    \item \textbf{Raspberry Pi:} Computador de placa única
        (\textit{Single-Board Computer}) com elevado poder computacional e suporte a
        sistema operacional Linux, amplamente utilizado em painéis digitais locais e
        processamento multimídia \cite{raspberrypi_datasheet}, porém com maior custo
        energético e financeiro, além de exigir controladores \acrshort{can} externos
        (como o circuito integrado MCP2515) via barramento SPI
        \cite{microchip_mcp2515}.
    \item \textbf{ESP32:} Sistema em um \textit{chip} (\textit{System-on-Chip}) de
        núcleo duplo da Espressif Systems (ilustrado na Figura~\ref{fig:esp32}) que se
        destaca pela conectividade \acrshort{wifi} e \textit{Bluetooth} integrada,
        suporte a aceleração criptográfica em \textit{hardware} e presença de um
        controlador periférico nativo compatível com a norma ISO 11898-1 denominado
        \acrshort{twai} (\textit{Two-Wire Automotive Interface}), exigindo apenas um
        transceptor diferencial externo de $3{,}3\text{ V}$ para acoplamento físico ao
        barramento \acrshort{can} veicular
        \cite{espressif_esp32_datasheet,ti_sn65hvd230}.
\end{itemize}

\begin{figure}[!htb]
	\centering
	\caption{Microcontrolador ESP32}
	\includegraphics[width=0.55\textwidth]{imagens/microcontrolador_esp32_vernizzi.jpeg}
	\legend{Fonte: Retirado de \citeonline{espressif_esp32_datasheet}.}
	\label{fig:esp32}
\end{figure}

\section{Sistemas Operacionais de Tempo Real (RTOS) e concorrência em sistemas
embarcados}\label{sec:rtos_concorrencia}

Em sistemas embarcados automotivos nos quais a aquisição de quadros de barramento em
alta frequência deve coexistir com operações de rede sem fio e criptografia, o modelo
sequencial de laço único (\textit{superloop}) mostra-se inadequado devido ao forte
acoplamento temporal entre rotinas \cite{ganssle2008,peckol2008}. Nesse modelo
monolítico sem preempção, o tempo de ciclo torna-se variável e dependente do pior
caminho de execução: o bloqueio de uma única rotina de entrada/saída ou de negociação de
rede paralisa todo o processador, ocasionando perda de mensagens do veículo.

Para superar as limitações do laço sequencial e compreender a arquitetura de
escalonamento adotada neste projeto, faz-se necessário distinguir um Sistema Operacional
de Propósito Geral (\acrshort{gpos}, do inglês
\textit{General-Purpose Operating System}, a exemplo de distribuições Linux e Windows)
de um Sistema Operacional de Tempo Real (\acrshort{rtos}, do inglês
\textit{Real-Time Operating System}). Segundo \citeonline{tanenbaum2006} e
\citeonline{peckol2008}, o escalonador de um \acrshort{gpos} é projetado para otimizar a
vazão média de processamento (\textit{throughput}) e a repartição justa de tempo
(\textit{fairness} ou \textit{time-sharing}) entre múltiplos processos interativos de
usuário, admitindo variações estocásticas de latência (\textit{jitter}) na ordem de
dezenas ou centenas de milissegundos. Em contrapartida, em um \acrshort{rtos}, a
correção do sistema não depende apenas da exatidão lógica ou aritmética do resultado
produzido, mas fundamentalmente do instante físico em que esse resultado é entregue em
relação a prazos estritos (\textit{deadlines}).

Os sistemas de tempo real classificam-se em rígidos (\textit{hard real-time}), nos quais
a violação de um prazo acarreta falha crítica ou perda irreversível de amostras de um
barramento físico, e brandos (\textit{soft real-time}), nos quais atrasos esporádicos
provocam apenas degradação tolerável de qualidade. Para assegurar esse determinismo
temporal, o núcleo (\textit{kernel}) de um \acrshort{rtos} implementa escalonamento
preemptivo baseado em prioridades fixas, garantindo que uma tarefa crítica desbloqueada
por um evento externo assuma imediatamente a unidade central de processamento em
detrimento de tarefas de menor prioridade, mantendo o tempo de troca de contexto
(\textit{context switch}) e a latência de interrupção estritamente limitados e
previsíveis \cite{ganssle2008,espressif2026}.

Todavia, a execução concorrente de múltiplas tarefas que compartilham recursos de
memória introduz o problema clássico do entrelaçamento não determinístico de instruções
(\textit{interleaving}) dentro de regiões críticas (\textit{critical sections}).
Conforme formalizado por \citeonline{benari2006}, a ausência de sincronização rigorosa
no acesso a variáveis compartilhadas resulta em condições de corrida
(\textit{race conditions}) e leituras inconsistentes (\textit{dirty reads}), ao passo
que o gerenciamento inadequado de bloqueios pode induzir anomalias severas de
escalonamento, tais como a inanição de tarefas de menor prioridade
(\textit{starvation}), o impasse circular (\textit{deadlock}) e a inversão de prioridade
(\textit{priority inversion}) \cite{tanenbaum2006}. A engenharia de \textit{firmware}
sob um \acrshort{rtos} exige, portanto, o uso disciplinado de primitivas formais de
exclusão mútua e comunicação intertarefas providas pelo \textit{kernel}.

\subsection{O \textit{framework} de desenvolvimento oficial ESP-IDF}\label{sec:espidf}

O \acrshort{espidf} (\textit{Espressif IoT Development Framework}) constitui o ambiente
oficial de desenvolvimento de \textit{software} em nível industrial mantido pela
Espressif Systems para a família de processadores ESP32 \cite{espressif2026}. Ao
contrário de plataformas de prototipagem rápida como o \textit{Arduino Core}, que
encapsulam a execução sob uma abstração cooperativa simplificada (\texttt{setup()} e
\texttt{loop()}) e ocultam o controle fino sobre registradores e interrupções, o
\acrshort{espidf} fornece acesso determinístico em linguagem C nativa à Camada de
Abstração de \textit{Hardware} (\acrshort{hal}, do inglês
\textit{Hardware Abstraction Layer}), aos periféricos de silício (como o controlador
automotivo \acrshort{twai} e os aceleradores criptográficos de \textit{hardware}) e às
pilhas de rede e segurança integradas, incluindo o protocolo TCP/IP leve
(\textit{LwIP}), a biblioteca criptográfica \textit{mbedTLS} e o cliente nativo
\texttt{esp-mqtt}.

Na engenharia de \textit{firmware}, o \acrshort{espidf} substitui arquivos de compilação
monolíticos por um sistema de meta-construção baseado em \textit{CMake} e no motor
\textit{Ninja}, parametrizado estaticamente em tempo de compilação via árvore
\textit{Kconfig} (\texttt{sdkconfig}). Essa infraestrutura viabiliza a adoção direta de
uma Arquitetura Baseada em Componentes (\acrshort{cba}, formalizada pela norma
ISO/IEC/IEEE 42010), na qual cada subsistema de \textit{software} é isolado em um
diretório autocontido, compilado pela cadeia de ferramentas cruzada
(\texttt{xtensa-esp32-elf-gcc}) como uma biblioteca estática independente
(\texttt{lib*.a}) e registrado com interfaces públicas explícitas
(\texttt{INCLUDE\_DIRS}) e grafos estritos de dependência (\texttt{REQUIRES})
\cite{iso42010}. Aliada às diretrizes de confiabilidade do padrão automotivo MISRA
C:2012 e do guia de codificação embarcada de \citeonline{barr2008}, essa estrutura
modular permite ainda embutir artefatos criptográficos diretamente na memória
\textit{Flash} durante a etapa de ligação (\textit{link-time embedding} via diretiva
\texttt{EMBED\_TXTFILES}), assegurando baixo acoplamento, rastreabilidade e segurança
operacional sem dependência de sistemas de arquivos externos.

\subsection{Integração arquitetural entre o ESP-IDF e o \textit{kernel}
FreeRTOS}\label{sec:integracao_freertos}

Diferentemente de microcontroladores convencionais nos quais um sistema operacional de
tempo real é adicionado opcionalmente como biblioteca externa sobre um código
\textit{bare-metal}, no ecossistema \acrshort{espidf} o \textit{kernel} FreeRTOS
constitui a fundação operacional nativa e indissociável do próprio \textit{framework}
\cite{espressif2026}. Enquanto a distribuição original do FreeRTOS foi concebida para
processadores de núcleo único, a Espressif implementou no \acrshort{espidf} uma extensão
arquitetural de Multiprocessamento Simétrico (\acrshort{smp}, do inglês
\textit{Symmetric Multiprocessing}), capaz de escalonar tarefas simultaneamente através
dos dois núcleos físicos de $32\text{ bits}$ Tensilica Xtensa LX6 do ESP32: o núcleo de
protocolo (\texttt{PRO\_CPU}, \textit{Core 0}), dedicado prioritariamente às tarefas
internas de rádio \acrshort{wifi}, pilha TCP/IP e eventos de sistema, e o núcleo de
aplicação (\texttt{APP\_CPU}, \textit{Core 1}), disponível para a execução paralela das
tarefas de usuário sob uma cadência de relógio de sistema (\textit{system tick}) de
$1000\text{ Hz}$ (fatias temporais determinísticas de $1\text{ ms}$). Durante a
inicialização do chip, o próprio código de partida do \acrshort{espidf} instancia os
serviços de base do \textit{kernel} (como o despachante de eventos \texttt{esp\_event} e
o serviço de temporizadores \texttt{Tmr Svc}) e lança a função de entrada
\texttt{app\_main()} já como uma tarefa nativa do FreeRTOS, a partir da qual os módulos
da aplicação criam suas próprias tarefas trabalhadoras com afinidade de núcleo definida
ou balanceada.

Essa simbiose profunda manifesta-se no acoplamento direto entre os \textit{drivers} de
periféricos do \acrshort{espidf} e as primitivas formais de sincronização e comunicação
intertarefas do FreeRTOS:
\begin{itemize}
    \item \textbf{Filas orientadas a interrupção (\textit{ISR-safe Queues}) nos
        \textit{drivers} de barramento:} Em consonância com os princípios de projeto de
        sistemas embarcados que preconizam rotinas de interrupção mínimas e velozes
        \cite{ganssle2008}, os \textit{drivers} nativos do \acrshort{espidf}, a exemplo
        do controlador automotivo \acrshort{twai}, operam sem espera ocupada
        (\textit{polling}). Quando um quadro \acrshort{can} é recebido no barramento, a
        rotina de serviço de interrupção (\acrshort{isr}) do \textit{driver} transfere o
        pacote diretamente para uma fila interna do FreeRTOS
        (\texttt{xQueueSendFromISR}) sob o modelo produtor-consumidor limitado
        (\textit{bounded buffer}) \cite{benari2006} e solicita ao escalonador o
        desbloqueio imediato da tarefa de recepção, garantindo latência mínima de
        resposta sem consumir ciclos de CPU em períodos de ociosidade.
    \item \textbf{Temporizadores de \textit{software} (\textit{Software Timers}) não
        bloqueantes:} A integração do serviço de temporizadores do FreeRTOS ao
        \acrshort{espidf} permite substituir atrasos bloqueantes (\texttt{vTaskDelay})
        por funções de retorno (\textit{callbacks}) assíncronas em segundo plano,
        evitando a inanição de processamento (\textit{starvation}) e prevenindo o
        transbordamento (\textit{overflow}) das filas de recepção \acrshort{twai}
        durante rajadas intensas de tráfego veicular \cite{ganssle2008,espressif2026}.
    \item \textbf{Grupos de eventos (\textit{Event Groups}) e semáforos de exclusão
        mútua (\textit{Mutex}):} A malha de eventos assíncronos do \acrshort{espidf}
        (responsável pelas transições de estado da conexão \acrshort{wifi} e concessão
        de endereço IP via DHCP) sincroniza-se deterministicamente com as tarefas de
        aplicação por meio de barreiras baseadas em \textit{Event Groups} do FreeRTOS.
        Simultaneamente, o compartilhamento de estruturas de telemetria em memória RAM
        entre as tarefas concorrentes do barramento \acrshort{can} e do cliente
        \acrshort{mqtts} é protegido por semáforos do tipo \textit{Mutex} dotados de
        mecanismo de herança de prioridade (\textit{priority inheritance}), assegurando
        a propriedade formal de exclusão mútua na seção crítica e erradicando condições
        de corrida (\textit{race conditions}), leituras corrompidas (\textit{dirty
        reads}) e inversões de prioridade \cite{benari2006,espressif2026}.
\end{itemize}

\section{Protocolo MQTT e segurança na camada de transporte (TLS) para
IoT}\label{sec:mqtt_tls}

A integração de dispositivos embarcados veiculares a plataformas de nuvem segue a cadeia
funcional típica de sistemas de \acrshort{iot}, compreendendo a coleta em sensores e
barramentos de borda, a transmissão por enlaces de comunicação sem fio, o processamento
e armazenamento em servidores remotos e a apresentação em interfaces de usuário,
conforme esquematizado na Figura~\ref{fig:sistema_iot} \cite{andrijauskas2018}. Contudo,
a transmissão contínua de séries temporais entre veículos em movimento e servidores em
nuvem impõe desafios severos relacionados à instabilidade de enlaces móveis, restrição
de largura de banda e segurança cibernética. Nesse contexto, o protocolo \acrshort{mqtt}
(\textit{Message Queuing Telemetry Transport}), padronizado pela organização OASIS e
pela norma ISO/IEC 20922, estabelece um modelo de mensageria leve fundamentado no
paradigma publicador/subscritor (\textit{publish/subscribe}) intermediado por um
servidor central (\textit{broker}) \cite{oasis2014}.

\begin{figure}[!htb]
	\centering
	\caption{Arquitetura funcional em blocos de um sistema de Internet das Coisas (IoT)}
	\includegraphics[width=0.78\textwidth]{imagens/sistema iot.jpg}
	\legend{Fonte: Adaptado de \citeonline{andrijauskas2018}.}
	\label{fig:sistema_iot}
\end{figure}

Diferentemente do protocolo \acrshort{http} tradicional, cuja semântica de requisição e
resposta síncrona cliente-servidor na camada de aplicação exige cabeçalhos textuais
extensos e repetidas negociações de conexão sobre o protocolo TCP \cite{kurose2010}, o
\acrshort{mqtt} opera sobre conexões TCP persistentes com cabeçalho fixo mínimo de
apenas $2\text{ bytes}$, promovendo o desacoplamento espacial, temporal e de
sincronização entre a unidade telemétrica embarcada (publicadora) e o serviço de
ingestão em nuvem (subscritor) \cite{iso20922}. A organização hierárquica em árvores de
tópicos possibilita ainda que o servidor de aplicação utilize caracteres coringa
(\textit{wildcards}, como o operador de nível único \texttt{+}) exclusivamente na
operação de subscrição para agregar o fluxo de múltiplos veículos de forma escalável.

\subsection{Vulnerabilidades do MQTT em texto claro e a adoção do MQTTS sobre
TLS}\label{sec:vulnerabilidades_mqtt}

Sob a ótica dos fundamentos de Segurança da Informação regidos pelas normas ISO/IEC
27001 e ISO/IEC 27002, a proteção de qualquer arquitetura de comunicação distribuída
estrutura-se sobre a preservação da tríade de Confidencialidade, Integridade e
Disponibilidade, somada à autenticidade das partes comunicantes \cite{baars2018}. No
contexto específico da Internet das Coisas aplicada a veículos conectados,
\citeonline{iqbal2019} ressaltam que a transmissão de dados telemétricos exige garantias
estritas de que as leituras provenham de uma fonte autêntica e íntegra, empregando
protocolos de comunicação dotados de mecanismos robustos contra o vazamento de
informações e tentativas de invasão ao longo de enlaces sem fio móveis (\acrshort{wifi}
e redes celulares 4G/LTE).

Embora o protocolo \acrshort{mqtt} padrão (operando na porta TCP 1883) apresente elevada
eficiência de transporte, sua especificação base não incorpora cifragem nativa na camada
de aplicação, trafegando credenciais de conexão, nomes de tópicos e cargas úteis
\acrshort{json} inteiramente em texto claro (\textit{plaintext}) \cite{oasis2014}. A
transmissão não criptografada em redes sem fio abertas ou compartilhadas expõe a
telemetria veicular a dois vetores clássicos de ameaça:

\begin{itemize}
    \item \textbf{Escuta clandestina ou interceptação passiva (\textit{Eavesdropping}):}
        Consiste na captura e inspeção não autorizada de pacotes de dados que transitam
        por um canal de comunicação por meio de analisadores de protocolo ou farejadores
        de rede (\textit{packet sniffers}) \cite{kurose2010,eleuterio2011}. Como o
        \acrshort{mqtt} puro transmite os segmentos TCP sem proteção criptográfica, um
        agente malicioso posicionado no mesmo enlace sem fio ou nos roteadores
        intermediários consegue reconstruir integralmente o conteúdo das mensagens
        \acrshort{json}, extraindo padrões operacionais do veículo e metadados da árvore
        de tópicos sem deixar rastros de modificação no fluxo de rede, o que viola
        diretamente o princípio da confidencialidade.
    \item \textbf{Ataque de intermediário (\textit{Man-in-the-Middle} ---
        \acrshort{mitm}):} Caracteriza um ataque ativo de rede no qual o invasor se
        interpõe secretamente no canal de comunicação entre o dispositivo embarcado
        (cliente) e o servidor (\textit{broker}), personificando cada extremidade
        legítima perante a outra \cite{baars2018,tanenbaum2011}. Diferentemente da
        escuta passiva, o atacante em posição de \acrshort{mitm} detém a capacidade de
        interceptar, descartar, reter ou adulterar ativamente os pacotes telemétricos em
        trânsito (por exemplo, injetando leituras falsas de rotação ou mascarando um
        superaquecimento crítico do motor), além de poder redirecionar a unidade
        embarcada para um servidor fraudulento (\textit{server spoofing}),
        comprometendo simultaneamente a confidencialidade, a integridade dos dados e a
        autenticidade da origem.
\end{itemize}

Para neutralizar esses vetores de ataque em conformidade com o paradigma de arquitetura
\textit{Zero Trust} (formalizado pelo NIST SP 800-207, segundo o qual nenhum enlace de
rede externa deve ser considerado implicitamente confiável) \cite{rose2020},
substitui-se o transporte em texto claro pelo encapsulamento da sessão \acrshort{mqtt}
sobre o protocolo criptográfico \acrshort{tls} (\textit{Transport Layer Security},
especificado pela RFC 5246), instituindo o padrão seguro \acrshort{mqtts} na porta 8883
\cite{dierks2008}. Durante o aperto de mão (\textit{handshake}) criptográfico do
\acrshort{tls}, o microcontrolador valida matematicamente a assinatura digital do
certificado X.509 apresentado pelo \textit{broker} perante a Autoridade Certificadora
(\texttt{CA}) raiz confiável gravada em sua memória \textit{Flash}, impedindo a
personificação de servidor típica de ataques \acrshort{mitm}. Em seguida,
estabelecem-se chaves de sessão simétricas efêmeras que cifram todo o fluxo de pacotes
(neutralizando o \textit{eavesdropping}) e anexam Códigos de Autenticação de Mensagem
(MAC) a cada registro transmitido, garantindo que qualquer tentativa de adulteração de
bytes em trânsito seja imediatamente detectada e descartada antes de atingir a camada de
persistência em nuvem.

\section{Mineração de dados (CRISP-DM) e detecção de anomalias via aprendizado de
máquina}\label{sec:crispdm_iforest}

A transformação de dados telemétricos brutos em diagnósticos preditivos requer uma
metodologia sistemática de ciência de dados. O modelo de referência \acrshort{crispdm}
(\textit{Cross-Industry Standard Process for Data Mining}), proposto por
\citeonline{wirth2000}, estrutura o ciclo de vida analítico em seis fases iterativas:
(i) Compreensão do Negócio (\textit{Business Understanding}), na qual se definem os
requisitos operacionais e os limites termomecânicos do motor; (ii) Compreensão dos Dados
(\textit{Data Understanding}), voltada à análise exploratória das séries temporais
coletadas; (iii) Preparação dos Dados (\textit{Data Preparation}), abrangendo limpeza,
tratamento de valores ausentes e engenharia de atributos; (iv) Modelagem
(\textit{Modeling}), etapa de estruturação dos algoritmos de aprendizado de máquina; (v)
Avaliação (\textit{Evaluation}), na qual se verifica a coerência física e estatística
das predições; e (vi) Implantação (\textit{Deployment}), integrando o modelo ao sistema
de monitoramento.

No domínio automotivo real, a obtenção de conjuntos de dados balanceados contendo falhas
mecânicas severas previamente rotuladas é inviável sem submeter o veículo a condições
destrutivas. Por essa razão, técnicas de aprendizado não supervisionado mostram-se
particularmente adequadas para a detecção de anomalias (\textit{outlier detection}).
Destaca-se nesse escopo o algoritmo \textit{Isolation Forest}, introduzido por
\citeonline{liu2008}, o qual identifica observações anômalas pelo princípio do
isolamento estocástico em vez de depender de métricas computacionalmente custosas de
distância ou densidade \cite{geron2022}.

O algoritmo \textit{Isolation Forest} constrói um conjunto (\textit{ensemble}) de
árvores binárias de isolamento (\textit{iTrees}) particionando recursivamente e de forma
aleatória o espaço de atributos (como pares de rotação em \acrshort{rpm}, velocidade e
temperatura do arrefecimento). Como pontos anômalos situam-se em regiões esparsas da
distribuição multivariada, eles exigem um número significativamente menor de partições
aleatórias para serem isolados nos nós folhas, resultando em comprimentos médios de
caminho (\textit{path lengths}) mais curtos desde a raiz da árvore \cite{liu2008}. Dessa
forma, o modelo fornece uma base matemática eficiente para detectar precocemente desvios
térmicos ou cinemáticos atípicos no funcionamento do motor com complexidade temporal
linear e baixo consumo de memória.

\chapter{Revisão da Literatura}

A literatura sobre diagnóstico e monitoramento veicular via porta \acrshort{obd}
apresenta diversas abordagens que variam em \textit{hardware}, \textit{software},
arquitetura de comunicação e objetivos analíticos. Esta revisão analisa um conjunto de
trabalhos correlatos para contextualizar e destacar as contribuições do presente
projeto. A análise está organizada por eixos temáticos, comparando as soluções
existentes com a arquitetura proposta neste trabalho quanto à concepção do dispositivo
de aquisição embarcado, plataformas de processamento e comunicação em nuvem, finalidade
analítica e escalabilidade de protocolos.

\section{\textit{Hardware}: \textit{scanners} comerciais e soluções customizadas}

A análise da literatura revela duas estratégias principais para a aquisição de dados da
rede veicular: o emprego de adaptadores diagnósticos comerciais prontos ou o
desenvolvimento de \textit{hardware} embarcado dedicado.

A primeira abordagem, focada na conveniência de integração, apoia-se em interfaces
comerciais fechadas. \citeonline{junior2009} exemplifica essa estratégia ao empregar um
cabo comercial VAG-COM para estabelecer comunicação com a \acrshort{ecu} de um veículo
Volkswagen Gol sob o protocolo proprietário KW1281. De maneira análoga,
\citeonline{andria2016}, no desenvolvimento de uma plataforma de aquisição veicular em
Bari, na Itália, e \citeonline{marques2019}, na construção do sistema
\textit{Vehicle-Manager}, utilizaram adaptadores comerciais baseados no circuito
integrado ELM327 com comunicação \textit{Bluetooth}. Revisões conduzidas por
\citeonline{manalu2017} e \citeonline{moniaga2018} também avaliaram diferentes
adaptadores comerciais \acrshort{obd} (\textit{iSaddle}, \textit{ERUSUN} e
\textit{iOBD2}) acoplados a computadores de placa única. Embora reduzam o tempo inicial
de prototipagem, tais dispositivos operam como ``caixas-pretas'' que traduzem comandos
seriais com latência adicional, restringindo o desenvolvedor às taxas de amostragem e
aos conjuntos de parâmetros impostos pelo fabricante do adaptador.

Em contrapartida, a segunda abordagem prioriza o controle sobre as camadas física e de
enlace por meio do projeto de \textit{hardware} customizado conectado diretamente ao
barramento \acrshort{can}. Nessa vertente, \citeonline{farinelli2016} projetou um
sistema embarcado de bordo dedicado à leitura direta de parâmetros de diagnóstico na
rede \acrshort{can}, detalhando o tratamento de quadros e a decodificação local de
\acrshort{pid}s. \citeonline{yu2024} desenvolveu um \textit{datalogger} automotivo
utilizando um computador de placa única Raspberry Pi acoplado via barramento SPI ao
controlador e transceptor \acrshort{can} MCP2515, provendo coleta e visualização em
tempo real por meio de uma Interface Humano-Computador (\acrshort{ihc}) local. Por sua
vez, \citeonline{vernizzi2024} implementou uma unidade de aquisição baseada no
microcontrolador ESP32 acoplado a um transceptor \acrshort{can} dedicado para leitura
direta da rede veicular, combinando transmissão sem fio via \acrshort{wifi} e
armazenamento redundante (\textit{backup}) em cartão de memória \acrshort{sd}.

Nesse contexto, o presente trabalho alinha-se à segunda vertente ao utilizar o
microcontrolador ESP32 integrado diretamente a um transceptor diferencial \acrshort{can}
de $3{,}3\text{ V}$ operando sobre o controlador periférico nativo \acrshort{twai}. Essa
decisão de engenharia elimina a dependência de adaptadores intermediários baseados em
ELM327 \cite{elmelectronics_elm327} e garante controle determinístico em tempo real
sobre a temporização das requisições \acrshort{sae} J1979 e a filtragem de interrupções
do barramento.

\section{Plataformas de \textit{software} e de comunicação}

A arquitetura de processamento, transporte e persistência dos dados coletados também
apresenta forte evolução na literatura. Diversas soluções concentram o processamento
exclusivamente no interior do veículo: o sistema de \citeonline{junior2009} exigia um
computador portátil (\textit{notebook}) executando \textit{software} em C\# como
intermediário obrigatório para ler a interface VAG-COM e retransmitir os dados via
\textit{Bluetooth}. Projetos mais recentes de computadores de bordo e painéis digitais,
como os desenvolvidos por \citeonline{cowfish2014}, \citeonline{wacker2017} (sistema
\textit{M3 Pi} com tela sensível ao toque e registro local em arquivos CSV) e
\citeonline{shinde2024} (painel digital em Python/PyQt5 sobre Raspberry Pi), eliminaram
o computador portátil, mas mantiveram os dados confinados ao dispositivo de visualização
local no painel do veículo.

Buscando superar o isolamento local, \citeonline{andria2016} implementaram em um
Raspberry Pi embarcado um banco de dados relacional MySQL local que armazena as séries
temporais durante o trajeto e realiza a sincronização remota via \acrshort{wifi} quando
o veículo retorna à garagem. Já \citeonline{marques2019} utilizou um aplicativo móvel
Android conectado via \textit{Bluetooth} ao veículo como ponte de transmissão para uma
interface \textit{web} centralizada. Reconhecendo a necessidade de conectividade
heterogênea diretamente a partir do conector diagnóstico, \citeonline{baek2015}
projetaram um conector \acrshort{obd} integrado baseado no microcontrolador ARM
Cortex-M3 (\texttt{STM32F103ZET6}) dotado simultaneamente de interfaces
\textit{Bluetooth}, \acrshort{wifi} e rede celular 3G WCDMA para transmissão direta a
dispositivos externos.

Mais recentemente, arquiteturas orientadas à \acrshort{iot} passaram a adotar protocolos
de mensageria assíncrona. \citeonline{oliveira2022} desenvolveu um sistema embarcado
combinando o microcontrolador ESP32 e a plataforma BeagleBone Black para transmitir
telemetria veicular em tempo real via protocolo \acrshort{mqtt} sobre rede celular
móvel, persistindo os registros em um banco de dados não relacional RavenDB na nuvem.
Avançando sobre essa abordagem, o presente trabalho elimina a necessidade de um segundo
computador embarcado intermediário (como a BeagleBone Black ou o Raspberry Pi) ao
executar toda a pilha criptográfica \acrshort{mqtts} (\acrshort{tls} v1.2) diretamente
no microcontrolador ESP32 sob o \textit{kernel} FreeRTOS, acoplando-o a um
\textit{backend} assíncrono com validação estrita de contratos e persistência relacional
transacional \acrshort{acid}.

\section{Finalidade da aplicação e profundidade da análise}

Os objetivos analíticos das plataformas reportadas na literatura podem ser classificados
em três níveis: monitoramento descritivo, diagnóstico reativo e análise comportamental
ou preditiva.

No nível reativo e descritivo, \citeonline{singh2021} propuseram um sistema de
diagnóstico baseado em \acrshort{iot} utilizando Raspberry Pi (biblioteca
\textit{pyOBD}), servidor em Node.js e banco de dados não relacional MongoDB com foco na
leitura e interpretação de códigos de falha (\acrshort{dtc}s) e envio de notificações de
manutenção ao usuário. No âmbito de gestão de frotas, \citeonline{marques2019}
direcionou sua plataforma ao acompanhamento histórico de quilometragem, temperatura e
alertas preventivos parametrizados.

Avançando para a análise cinemática, \citeonline{andria2016} integraram aos dados de
rotação e velocidade do \acrshort{obd} uma Unidade de Medição Inercial (\acrshort{imu}
MPU-6050 via barramento I2C) e um receptor GPS para avaliar o comportamento do motorista
durante as trocas de marcha por meio da razão instantânea velocidade/\acrshort{rpm},
classificando o estilo de condução visando à segurança e à economia de combustível. De
forma complementar, \citeonline{oliveira2022} utilizou leituras de acelerometria e
velocidade para a detecção determinística de eventos de frenagem brusca.

Posicionando-se além da leitura reativa de códigos \acrshort{dtc} \cite{singh2021} ou da
aplicação de limiares estáticos sobre grandezas isoladas
\cite{marques2019,oliveira2022}, a arquitetura proposta neste trabalho estrutura a
consolidação de séries temporais multivariadas do motor (rotação, velocidade e
temperatura do líquido de arrefecimento) para habilitar uma camada analítica superior: a
modelagem de detecção não supervisionada de anomalias baseada em aprendizado de máquina
(\textit{Isolation Forest}) sob a metodologia \acrshort{crispdm}, visando identificar
padrões sutis de degradação termomecânica antes que um código de defeito seja disparado
pela \acrshort{ecu}.

\section{Escalabilidade e flexibilidade de protocolos}

A escalabilidade de sistemas de telemetria automotiva depende diretamente do grau de
desacoplamento entre o \textit{hardware} de borda e a infraestrutura de nuvem. Soluções
fortemente acopladas a protocolos legados proprietários, como a implementação de
\citeonline{junior2009} restrita ao padrão KW1281, apresentam baixa portabilidade entre
diferentes fabricantes e modelos de veículos.

Na arquitetura desenvolvida neste trabalho, a escalabilidade é assegurada em múltiplas
camadas. Na borda, a adoção do padrão internacional ISO 15765-4 (\acrshort{can} a
$500\text{ kbaud}$) e a estruturação do \textit{firmware} em C sob Arquitetura Baseada
em Componentes (\acrshort{cba}) permitem que novos identificadores \acrshort{pid} da
norma \acrshort{sae} J1979 sejam incorporados sem impactar os módulos de conectividade e
segurança. Na nuvem, o modelo publicador/subscritor do protocolo \acrshort{mqtts} com
subscrição via caractere coringa (\textit{wildcard}) permite que um único serviço de
ingestão processe simultaneamente o fluxo telemétrico de múltiplos veículos
(\textit{multi-tenant}), garantindo flexibilidade para a evolução contínua dos modelos
analíticos.

\section{Considerações finais do capítulo}

A partir da análise comparativa da literatura, constata-se que o presente trabalho se
diferencia por integrar quatro eixos complementares em uma arquitetura coesa: (i)
aquisição embarcada nativa em barramento \acrshort{can} via periférico \acrshort{twai}
do ESP32 sob \acrshort{rtos}, superando as limitações de adaptadores comerciais baseados
em ELM327 \cite{andria2016,marques2019}; (ii) transmissão direta da borda para a nuvem
sob canal criptografado \acrshort{mqtts}/\acrshort{tls}, dispensando computadores
intermediários no interior do veículo \cite{junior2009,oliveira2022}; (iii) ingestão
assíncrona com validação formal de contratos e persistência relacional \acrshort{acid}
aliada à visualização em tempo real; e (iv) modelagem analítica voltada à detecção não
supervisionada de anomalias operacionais, avançando em relação às abordagens puramente
descritivas ou reativas \cite{singh2021,yu2024}.

\chapter{Metodologia}

Para materializar os objetivos específicos traçados no Capítulo 1 e aplicar os
fundamentos teóricos discutidos no Capítulo 2 sobre as lacunas identificadas no Capítulo
3, este capítulo descreve as decisões de projeto e os procedimentos metodológicos
adotados na construção e validação da plataforma de telemetria automotiva inteligente. A
exposição organiza-se em seis frentes sequenciais e interdependentes: (i) evolução de
escopo, especificação de materiais e arquitetura geral do sistema; (ii) organização
modular do repositório de código-fonte no GitHub; (iii) engenharia de \textit{firmware}
embarcado em tempo real; (iv) engenharia de \textit{backend}, segurança cibernética e
persistência transacional; (v) emulação sintética veicular, visualização em tempo real e
modelagem matemática da camada de detecção de anomalias; e (vi) protocolo metodológico
de verificação de lacunas e revisão com apoio instrumental de Inteligência Artificial
Generativa (\acrshort{iag}).

\section{Evolução de escopo, materiais e visão geral da arquitetura}

Na concepção preliminar do projeto elaborada durante o projeto, o estudo de
caso previa a utilização de um veículo Volkswagen Polo Classic 2000, cuja interface
diagnóstica baseava-se no protocolo serial assíncrono legado \textit{K-Line}, regido
pela norma ISO 9141-2 a uma taxa de transmissão de $10{,}4\text{ kbaud}$ no pino 7 do
conector \acrshort{obd}, exigindo rotina temporizada de inicialização
(\textit{5-Baud Init}) e condicionamento de nível lógico por meio do transceptor
monolítico L9637D \cite{l9637d_datasheet}. Contudo, no início da execução experimental
do \acrshort{tcc} II, promoveu-se a evolução tecnológica do escopo para um veículo
contemporâneo, o Volkswagen Polo Track, equipado com barramento diferencial
\acrshort{can} operando a $500\text{ kbaud}$ sob a norma ISO 15765-4 \cite{iso15765}.

Em consonância com a convergência normativa e o ganho nominal de $48$ vezes em largura
de banda fundamentados na Seção~\ref{sec:convergencia_osi}, essa decisão de engenharia
eliminou os atrasos de sincronização inicial e assegurou aplicabilidade direta da
plataforma à frota automotiva moderna. Adicionalmente, a migração para o barramento
\acrshort{can} permitiu explorar a integração física direta com o controlador periférico
\acrshort{twai} presente no silício do microcontrolador ESP32, conforme caracterizado na
Seção~\ref{sec:microcontroladores} \cite{espressif_esp32_datasheet}. Sob essa
especificação, a arquitetura do sistema foi particionada em três subsistemas
desacoplados:

Sob essa nova especificação, a arquitetura do sistema foi projetada segundo os
princípios de engenharia de sistemas desacoplados, particionando o fluxo de
processamento desde a borda veicular até a nuvem em três subsistemas interdependentes:

\begin{itemize}
    \item \textbf{Subsistema de Borda Automotiva (\textit{Edge}):} Composto pela placa
        de desenvolvimento ESP32 DevKit V1 (equipada com o módulo ESP32-WROOM-32
        descrito
        na Seção~\ref{sec:microcontroladores}) acoplada ao transceptor diferencial
        \acrshort{can} SN65HVD230 da Texas Instruments \cite{ti_sn65hvd230}. A seleção
        do
        circuito integrado (\acrshort{ci}) SN65HVD230 justifica-se por sua operação
        nativa em tensão lógica de $3{,}3\text{ V}$, diretamente compatível com os pinos
        de entrada e saída de propósito geral (\acrshort{gpio}) do ESP32 sem a
        necessidade de conversores bidirecionais de nível lógico exigidos por
        transceptores convencionais de $5\text{ V}$. A alimentação elétrica do conjunto
        é
        derivada da linha automotiva de $12\text{ V}$ contínua (pino 16 do conector
        diagnóstico \acrshort{sae} J1962), regulada para $5\text{ V}$ no barramento de
        entrada da placa e compartilhando referência de terra comum (\texttt{GND}) com
        os pinos 4 e 5 do conector veicular.
    \item \textbf{Subsistema de Transporte Seguro e Ingestão em Nuvem (\textit{Cloud
        Backend}):} A unidade embarcada estabelece conectividade sem fio via enlace
        \acrshort{wifi} IEEE 802.11 b/g/n roteado por ponto de acesso móvel veicular
        4G/LTE, transmitindo cargas úteis em formato \acrshort{json} sobre o protocolo
        \acrshort{mqtts} (\acrshort{tls} v1.2 na porta 8883) para o \textit{broker}
        Eclipse Mosquitto. Na camada de servidor, um serviço de aplicação assíncrono
        desenvolvido em Python 3.12 sobre o \textit{framework} FastAPI e a biblioteca
        Eclipse Paho-MQTT executa a validação sintática e semântica de fronteira via
        Pydantic v2 e realiza a persistência relacional transacional em banco de dados
        SQLite, mediada pelo mapeamento objeto-relacional (\acrshort{orm}) SQLAlchemy
        2.0.
    \item \textbf{Subsistema de Visualização e Modelagem Preditiva:} Constituído pelo
        \textit{dashboard} analítico interativo desenvolvido em Streamlit, apoiado pelo
        emulador sintético de dinâmica veicular para ensaios de estresse e pela
        especificação arquitetural e modelagem matemática do \textit{pipeline} de
        detecção não supervisionada de anomalias estruturado sob a metodologia
        \acrshort{crispdm} e o algoritmo \textit{Isolation Forest}.
\end{itemize}

A topologia geral de acoplamento físico entre o barramento diferencial veicular
(\texttt{CAN\_H} e \texttt{CAN\_L}), o transceptor \acrshort{can} externo e as linhas
lógicas de transmissão (\texttt{CAN\_TX}) e recepção (\texttt{CAN\_RX}) roteadas pelos
pinos \acrshort{gpio} e barramentos de alimentação (\texttt{3V3}, \texttt{5V0} e
\texttt{GND}) do módulo ESP32-WROOM é ilustrada na Figura~\ref{fig:circuito_esp32}.
Complementarmente, o mapeamento elétrico definitivo adotado na implementação do
protótipo, incluindo a atribuição exata dos pinos \texttt{GPIO 4} e \texttt{GPIO 5}
para o controlador \acrshort{twai} e dos pinos de diagnóstico luminoso, encontra-se
consolidado no Quadro~\ref{quad:pinout}.

\begin{figure}[!htb]
	\centering
	\caption{Topologia de interfaceamento entre o barramento diferencial CAN, o transceptor
	externo e o microcontrolador ESP32}
	\includegraphics[width=0.78\textwidth]{imagens/circuito esp32.png}
	\legend{Fonte: Elaborado pelo autor (2026).}
	\label{fig:circuito_esp32}
\end{figure}

\begin{table}[!htb]
\centering
\refstepcounter{quadro}\label{quad:pinout}
{\footnotesize
\textbf{Quadro \thequadro{} -- Mapeamento elétrico de pinos entre o conector
\acrshort{obd} (\acrshort{sae} J1962), o transceptor SN65HVD230 e o microcontrolador
ESP32}\par}
\vspace{4pt}
{\footnotesize
\begin{tabular}{|>{\raggedright\arraybackslash}p{3.1cm}|l|>{\raggedright\arraybackslash}p{8.6cm}|}
\hline
\textbf{Pino / Barramento} & \textbf{Direção / Tipo}
    & \textbf{Função no Sistema e Conexão Física} \\ \hline
\texttt{GPIO 4} (\texttt{D4}) & Saída (\texttt{TX})
    & Linha de transmissão lógica \acrshort{twai}/\acrshort{can} conectada ao pino
    \texttt{CTX} (\texttt{TXD}) do transceptor SN65HVD230. \\ \hline
\texttt{GPIO 5} (\texttt{D5}) & Entrada (\texttt{RX})
    & Linha de recepção lógica \acrshort{twai}/\acrshort{can} conectada ao pino
    \texttt{CRX} (\texttt{RXD}) do transceptor SN65HVD230. \\ \hline
\texttt{GPIO 13} (\texttt{D13}) & Saída Digital
    & Indicador luminoso (LED) de sistema energizado e controlador \acrshort{twai}
    inicializado. \\ \hline
\texttt{GPIO 14} (\texttt{D14}) & Saída Digital
    & Indicador luminoso (LED) de tráfego \acrshort{can} válido recebido da
    \acrshort{ecu} (\texttt{0x7E8}). \\ \hline
\texttt{GPIO 26} (\texttt{D26}) & Saída Digital
    & Indicador luminoso (LED) de pacote telemétrico publicado com sucesso via
    \acrshort{mqtts}. \\ \hline
\texttt{3V3} e \texttt{GND} & Alimentação
    & Alimentação regulada de $3{,}3\text{ V}$ do transceptor SN65HVD230 e referência
    comum com os pinos 4 e 5 do conector \acrshort{obd}. \\ \hline
\texttt{CAN\_H} (Pino 6) e \texttt{CAN\_L} (Pino 14) & Diferencial
    & Barramento diferencial automotivo ISO 15765-4 do conector \acrshort{obd} conectado
    aos pinos \texttt{CANH} e \texttt{CANL} do SN65HVD230. \\ \hline
\end{tabular}\par}
\vspace{4pt}
{\footnotesize\raggedright
Fonte: Elaborado pelo autor (2026).\par}
\end{table}

\section{Organização do repositório e estrutura de diretórios no GitHub}

Para assegurar rastreabilidade, reprodutibilidade experimental e isolamento entre os
domínios de \textit{hardware} embarcado, infraestrutura de nuvem e visualização
analítica, todo o código-fonte do projeto foi consolidado em um repositório único
(\textit{monorepo}) versionado sob o sistema Git e hospedado na plataforma GitHub. A
árvore de diretórios foi projetada em estrita observância ao princípio de Separação de
Preocupações (\textit{Separation of Concerns}) \cite{martin2017}, concentrando os
artefatos de implementação em três núcleos funcionais complementados pelos diretórios de
segurança e documentação:

\begin{enumerate}
    \item \textbf{Núcleo de \textit{Hardware} e \textit{Firmware} Embarcado
        (\texttt{hardware/}):} Subdivide-se entre o diretório de prototipagem preliminar
        (\texttt{hardware/\allowbreak arduino\_sketch/}, apoiado pelo código de ensaio
        criptográfico \texttt{test/\allowbreak arduino/\allowbreak
        mqtts\_test/\allowbreak mqtts\_test.ino}) e o diretório de produção
        (\texttt{hardware/\allowbreak esp32/}) regido pelo \textit{framework} ESP-IDF.
        Neste último, a compilação global é parametrizada pelo manifesto raiz
        \texttt{CMakeLists.txt} e o ponto de entrada da aplicação reside em
        \texttt{main/\allowbreak main.c}. Os subsistemas de baixo nível encontram-se
        encapsulados em quatro pacotes independentes no diretório
        \texttt{hardware/\allowbreak esp32/\allowbreak components/}: \texttt{telemetry/}
        (gerenciamento concorrente de estado em memória), \texttt{obd2\_can/}
        (comunicação automotiva \acrshort{twai}/\acrshort{can}), \texttt{wifi\_manager/}
        (conectividade sem fio e sincronização de rede) e \texttt{mqtt\_service/}
        (transporte seguro \acrshort{mqtts} e certificado raiz embutido
        \texttt{ca\_cert.h}). Cada componente possui cabeçalho público próprio
        (\texttt{include/*.h}), unidade de tradução privada (\texttt{*.c}) e descritor
        de compilação local (\texttt{CMakeLists.txt}).
    \item \textbf{Núcleo de \textit{Backend}, Persistência e Testes
        (\texttt{backend/}):} Estruturado modularmente para desacoplar configuração,
        transporte assíncrono, validação e banco de dados. Compreende o módulo de
        parametrização centralizada (\texttt{backend/\allowbreak config.py}), a camada
        de persistência relacional e definição da entidade \acrshort{orm}
        (\texttt{backend/\allowbreak database.py} e
        \texttt{backend/\allowbreak models/\allowbreak telemetry.py}, responsáveis pelo
        gerenciamento do arquivo de banco de dados \texttt{vehicle\_data.db}), os
        contratos de validação de fronteira
        (\texttt{backend/\allowbreak schemas/\allowbreak telemetry.py}), o trabalhador
        assíncrono de subscrição de mensageria
        (\texttt{backend/\allowbreak mqtt\_client.py}) e o servidor de aplicação
        \textit{web} RESTful (\texttt{backend/\allowbreak main.py}). A verificação
        determinística dessa infraestrutura concentra-se no subdiretório
        \texttt{backend/\allowbreak tests/}, composto pelas suítes automatizadas
        \texttt{test\_main.py} e \texttt{test\_mqtt.py}.
    \item \textbf{Núcleo de Visualização, Emulação e Documentação (\texttt{frontend/},
        \texttt{scripts/}, \texttt{certs/} e \texttt{docs/}):} A interface gráfica de
        monitoramento em tempo real está implementada em
        \texttt{frontend/\allowbreak app.py}, enquanto o gerador sintético de tráfego
        telemétrico veicular reside em \texttt{scripts/\allowbreak mock\_car.py}. A
        cadeia de certificação pública X.509 da autoridade certificadora do
        \textit{broker} (\texttt{mosquitto.org.crt}) é mantida de forma centralizada
        para validação cruzada entre emulador e servidor, ao passo que o diretório
        \texttt{docs/} armazena a documentação arquitetural e os códigos-fonte
        tipográficos da presente monografia.
\end{enumerate}

O Quadro~\ref{quad:repo_estrutura} sintetiza o mapeamento completo dos diretórios e
artefatos de código que compõem o repositório do projeto, relacionando sua linguagem ou
tecnologia base à respectiva responsabilidade arquitetural.

\begin{table}[!htb]
\centering
\refstepcounter{quadro}\label{quad:repo_estrutura}
{\footnotesize
\textbf{Quadro \thequadro{} -- Estrutura modular de diretórios e artefatos de código no
repositório GitHub}\par}
\vspace{4pt}
{\footnotesize
\begin{tabular}{|>{\raggedright\arraybackslash}p{4.3cm}|>{\raggedright\arraybackslash}p{2.7cm}|>{\raggedright\arraybackslash}p{6.8cm}|}
\hline
\textbf{Diretório / Artefato} & \textbf{Tecnologia Base}
    & \textbf{Responsabilidade Arquitetural no Sistema} \\ \hline
\texttt{hardware/\allowbreak esp32/\allowbreak main/main.c} & C (ESP-IDF v6.1)
    & Ponto de entrada (\texttt{app\_main}) e orquestração da inicialização
    determinística dos componentes embarcados. \\ \hline
\texttt{hardware/\allowbreak esp32/\allowbreak components/\allowbreak telemetry/}
    & C / FreeRTOS
    & Repositório de estado telemétrico em memória volátil protegido por semáforo de
    exclusão mútua (\textit{Mutex}). \\ \hline
\texttt{hardware/\allowbreak esp32/\allowbreak components/\allowbreak obd2\_can/}
    & C / Driver \acrshort{twai}
    & Aquisição ativa de \acrshort{pid}s no barramento \acrshort{can}
    ($500\text{ kbaud}$) e conversão física segundo a norma \acrshort{sae} J1979. \\
    \hline
\texttt{hardware/\allowbreak esp32/\allowbreak components/\allowbreak wifi\_manager/}
    & C / ESP-Wi-Fi
    & Gerenciamento de conexão IEEE 802.11, reconexão automática e barreira de
    sincronização via \textit{EventGroup}. \\ \hline
\texttt{hardware/\allowbreak esp32/\allowbreak components/\allowbreak mqtt\_service/}
    & C / \texttt{esp-mqtt} (\acrshort{tls})
    & Publicação periódica de pacotes \acrshort{json} via \acrshort{mqtts} (porta 8883)
    com certificado X.509 embutido em memória \textit{Flash}. \\ \hline
\texttt{backend/\allowbreak config.py} & Python 3.12
    & Fonte Única da Verdade (\acrshort{ssot}) para constantes de rede, tópicos
    \acrshort{mqtt} e resolução de caminhos de arquivos. \\ \hline
\texttt{backend/\allowbreak database.py} e \texttt{models/\allowbreak telemetry.py}
    & SQLAlchemy 2.0 / SQLite
    & Configuração do motor relacional, controle de sessões transacionais
    \acrshort{acid} e mapeamento da tabela \texttt{telemetry}. \\ \hline
\texttt{backend/\allowbreak schemas/\allowbreak telemetry.py} & Pydantic v2
    & Validação sintática e semântica estrita de tipos e contratos de entrada e saída da
    \acrshort{api}. \\ \hline
\texttt{backend/\allowbreak mqtt\_client.py} & Eclipse Paho-MQTT
    & Trabalhador assíncrono de subscrição coringa (\textit{multi-tenant}) sobre
    \acrshort{tls} v1.2 e persistência defensiva. \\ \hline
\texttt{backend/\allowbreak main.py} & FastAPI / Uvicorn
    & Exposição de rotas RESTful (\texttt{GET}/\texttt{POST}) e gerenciamento do ciclo
    de vida (\texttt{lifespan}) do serviço de ingestão. \\ \hline
\texttt{backend/\allowbreak tests/\allowbreak test\_main.py} e \texttt{test\_mqtt.py}
    & \texttt{pytest} / \texttt{unittest.mock}
    & Suíte automatizada de testes unitários e de integração das rotas \acrshort{http} e
    do consumidor \acrshort{mqtts}. \\ \hline
\texttt{scripts/\allowbreak mock\_car.py} & Python 3.12 (\acrshort{mqtts})
    & Emulador sintético de dinâmica veicular para testes de carga e validação contínua
    sem dependência do veículo. \\ \hline
\texttt{frontend/\allowbreak app.py} & Streamlit / Pandas
    & \textit{Dashboard} analítico em tempo real com atualização seletiva sem
    \textit{flicker} e cálculo de grandezas cinemáticas. \\ \hline
\texttt{certs/\allowbreak mosquitto.org.crt} & Certificado X.509
    & Cadeia de autoridade certificadora raiz para autenticação criptográfica
    \acrshort{tls} v1.2 perante o \textit{broker} em nuvem. \\ \hline
\end{tabular}\par}
\vspace{4pt}
{\footnotesize\raggedright
Fonte: Elaborado pelo autor (2026).\par}
\end{table}

\section{Engenharia de \textit{firmware} embarcado}

Previamente à implementação definitiva do \textit{software} de baixo nível, conduziu-se
uma prova de conceito exploratória (\textit{Spike} tecnológico) utilizando o núcleo
Arduino (\textit{Arduino Core}) com as bibliotecas \texttt{WiFiClientSecure} e
\texttt{PubSubClient}, visando aferir a viabilidade do aperto de mão
(\textit{handshake}) criptográfico \acrshort{tls} v1.2 na porta 8883 com o certificado
raiz X.509 embutido em memória \textit{Flash}. Embora o experimento tenha comprovado a
capacidade da aceleração criptográfica por \textit{hardware} do ESP32 para processar
cifras assimétricas RSA e simétricas AES-GCM sem degradação severa de latência,
verificou-se empiricamente a limitação estrutural apontada na Seção~\ref{sec:espidf}: a
abstração cooperativa do ambiente Arduino restringia o controle determinístico sobre
prioridades de interrupção do \acrshort{rtos}, o acesso direto às filas de recepção do
periférico \acrshort{twai} e o isolamento estrito de memória. Por conseguinte, migrou-se
o desenvolvimento final para linguagem C nativa sobre o \textit{framework} oficial
\acrshort{espidf} v6.1 e o \textit{kernel} FreeRTOS em arquitetura \acrshort{smp},
compilado via sistema CMake/Ninja com a cadeia \texttt{xtensa-esp32-elf-gcc}
\cite{espressif2026}.

Aplicando o modelo de Arquitetura Baseada em Componentes (\acrshort{cba}) fundamentado
na Seção~\ref{sec:espidf} \cite{iso42010}, o código-fonte monolítico inicial foi
particionado em quatro bibliotecas estáticas independentes registradas via macro
\texttt{idf\_component\_register(...)}, coordenadas por um módulo orquestrador enxuto
cuja função principal \texttt{app\_main()} foi reduzida a 35 linhas dedicadas à
inicialização sequencial e verificação estrita de falhas (\texttt{ESP\_ERROR\_CHECK}).

O
primeiro componente implementado atua como o repositório central de estado telemétrico
em memória volátil por meio da estrutura \texttt{telemetry\_\allowbreak data\_t}, que
armazena as grandezas \texttt{rpm}, \texttt{temperature}, \texttt{speed}, a marca
temporal (\texttt{timestamp\_ms}) e sinalizadores booleanos de validade
(\texttt{rpm\_valid}, \texttt{temp\_\allowbreak valid} e \texttt{speed\_valid}). Em
aderência às diretrizes de segurança de código MISRA C:2012 \cite{misra2013} e aos
mecanismos de sincronização de seções críticas descritos na
Seção~\ref{sec:rtos_concorrencia}, a instância de dados é encapsulada no escopo restrito
do módulo com o qualificador \texttt{static}, e todas as rotinas de escrita atômica
(\texttt{telemetry\_set\_*}) e de extração de cópia consistente
(\texttt{telemetry\_\allowbreak get\_snapshot}) são protegidas por um semáforo de
exclusão mútua (FreeRTOS \textit{Mutex}, \texttt{Semaphore\allowbreak Handle\_t
s\_mutex}, com tempo limite de $100\text{ ms}$), impedindo condições de corrida e
leituras parciais entre tarefas concorrentes.

O segundo componente materializa a pilha de comunicação diagnóstica automotiva
estratificada na Seção~\ref{sec:convergencia_osi}, operando nas camadas física, de
enlace e de transporte da norma ISO 15765-4 \cite{iso15765} por meio do controlador
nativo \acrshort{twai} configurado a $500\text{ kbaud}$ via macro
\texttt{TWAI\_TIMING\_\allowbreak CONFIG\_\allowbreak 500KBITS()}, modo operacional
normal (\texttt{TWAI\_\allowbreak MODE\_\allowbreak NORMAL}) e filtro de aceitação
aberto (\texttt{TWAI\_\allowbreak FILTER\_\allowbreak CONFIG\_\allowbreak
ACCEPT\_\allowbreak ALL()}). A aquisição ativa de telemetria opera por meio de duas
tarefas concorrentes no FreeRTOS:

\begin{itemize}
    \item \textbf{Tarefa de transmissão (\texttt{twai\_transmit\_task}):} Executa em
        ciclo determinístico de $100\text{ ms}$, despachando quadros \acrshort{can}
        padrão de $11\text{ bits}$ com código de comprimento de dados (\acrshort{dlc})
        igual a $8\text{ bytes}$ direcionados ao identificador funcional de
        \textit{broadcast} \texttt{0x7DF}. A carga útil requisita no Modo de Serviço
        \texttt{0x01} os identificadores \acrshort{pid} \texttt{0x0C} (rotação do
        motor), \texttt{0x0D} (velocidade veicular) e \texttt{0x05} (temperatura do
        líquido de arrefecimento), preenchendo os bytes remanescentes com valores de
        \textit{padding} (\texttt{0x55} e \texttt{0xAA}).
    \item \textbf{Tarefa de recepção (\texttt{twai\_receive\_task}):} Consome de forma
        contínua os pacotes enfileirados pela rotina de serviço de interrupção
        (\acrshort{isr}) do \textit{driver} \acrshort{twai}, filtrando quadros
        originados pela \acrshort{ecu} principal do motor (\texttt{0x7E8}) que
        apresentem byte de confirmação de serviço \texttt{0x41}
        ($\texttt{0x01} + \texttt{0x40}$). Previamente à indexação dos vetores de carga
        útil, aplica-se validação defensiva do comprimento de dados
        ($\text{\acrshort{dlc}} \ge 5$ para o \acrshort{pid} de dois bytes e
        $\text{\acrshort{dlc}} \ge 3$ para grandezas de um byte), prevenindo falhas de
        segmentação de memória por quadros truncados ou ruidosos.
\end{itemize}

Na camada de aplicação diagnóstica (Seção~\ref{sec:aquisicao_obd2}), a conversão dos
bytes brutos hexadecimais retornados pela \acrshort{ecu}, nos quais $A$ e $B$
representam os bytes decimais de dados ($0$ a $255$), para grandezas físicas de
engenharia segue as expressões normativas definidas pela especificação \acrshort{sae}
J1979 \cite{saej1979}, implementadas conforme as Equações~\eqref{eq:rpm},
\eqref{eq:speed} e \eqref{eq:temp}:

\begin{equation}\label{eq:rpm}
\text{RPM} = \frac{256 \cdot A + B}{4}
\end{equation}

\begin{equation}\label{eq:speed}
v = A
\end{equation}

\begin{equation}\label{eq:temp}
T_{\text{ect}} = A - 40
\end{equation}

\noindent em que $\text{RPM}$ corresponde à rotação angular do eixo de manivelas (em
$\text{rpm}$, \acrshort{pid} \texttt{0x0C} na Equação~\eqref{eq:rpm}), $v$ representa a
velocidade instantânea de deslocamento do veículo (em $\text{km/h}$, \acrshort{pid}
\texttt{0x0D} na Equação~\eqref{eq:speed}) e $T_{\text{ect}}$ denota a temperatura do
líquido de arrefecimento do motor (em $^{\circ}\text{C}$, \acrshort{pid} \texttt{0x05}
na Equação~\eqref{eq:temp}, com deslocamento térmico de $-40\text{ }^{\circ}\text{C}$
para representação de temperaturas negativas em aritmética inteira sem sinal).
Adicionalmente, para erradicar gargalos temporais na recepção, substituiu-se o uso de
atrasos bloqueantes (\texttt{vTaskDelay}) no acionamento do indicador visual
\texttt{GPIO 14} por um temporizador de \textit{software} assíncrono do FreeRTOS
(\texttt{can\_led\_timer}, parametrizado em $30\text{ ms}$), liberando imediatamente a
tarefa de recepção para drenar a fila do controlador \acrshort{twai} sem perda de
quadros.

Por sua vez, o terceiro componente, encarregado da gerência de conectividade sem fio,
configura a interface de rádio em modo \textit{Station} (\texttt{WIFI\_MODE\_STA})
dotada de máquina de estados de reconexão automática
(\texttt{WIFI\_\allowbreak EVENT\_\allowbreak STA\_\allowbreak DISCONNECTED}) e barreira
de sincronização determinística via FreeRTOS \textit{EventGroup}
(\texttt{WIFI\_\allowbreak CONNECTED\_\allowbreak BIT}), a qual retém a inicialização da
camada de mensageria até a confirmação do evento
\texttt{IP\_\allowbreak EVENT\_\allowbreak STA\_\allowbreak GOT\_IP} pelo servidor DHCP
do roteador móvel. Finalmente, o quarto componente implementa o cliente de mensageria
segura baseado na biblioteca nativa \texttt{esp-mqtt}, apontado para o endereço
\texttt{mqtts://test.mosquitto.org:8883} sobre \acrshort{tls} v1.2 (pilha mbedTLS)
\cite{dierks2008}. O certificado raiz X.509 da autoridade certificadora é compilado
diretamente na imagem binária em memória \textit{Flash} via diretiva
\texttt{EMBED\_\allowbreak TXTFILES}, enquanto a tarefa periódica
\texttt{mqtt\_\allowbreak publisher\_\allowbreak task} opera em cadência de
$1000\text{ ms}$, adquire o \textit{snapshot} protegido por \textit{Mutex}, serializa a
carga útil em formato \acrshort{json}
(\texttt{\{"rpm": ..., "temperature": ..., "speed": ...\}}) e publica no tópico do
veículo, acionando o pino \texttt{GPIO 26} com desarme não bloqueante mediado por
temporizador FreeRTOS (\texttt{mqtt\_\allowbreak led\_timer}, $50\text{ ms}$). Uma vez
transmitidos os pacotes criptografados pela unidade embarcada, o fluxo de dados alcança
a infraestrutura de ingestão em nuvem.

\section{Engenharia de \textit{backend}, segurança e persistência}

O desenvolvimento da infraestrutura de servidor iniciou-se sob a estratégia de
engenharia \textit{Mock-First}, validando os contratos de dados e a persistência
relacional mediante requisições síncronas (\acrshort{http} POST), para em seguida
efetuar a transição arquitetural definitiva ao paradigma assíncrono orientado a eventos
via protocolo \acrshort{mqtts} na porta 8883 \cite{oasis2014,iso20922}. Em alinhamento
aos fundamentos de mensageria expostos na Seção~\ref{sec:mqtt_tls}, a árvore de tópicos
foi modelada em estrita observância à Seção 4.7 da especificação OASIS \acrshort{mqtt}
v3.1.1, que veda o uso de caracteres coringa (\textit{wildcards}) em operações de
publicação. Dessa forma, enquanto o microcontrolador ESP32 e o emulador publicam
(\texttt{PUBLISH}) exclusivamente em tópicos concretos e unívocos
(\texttt{ls-praxedes/\allowbreak scanner-obd2/\allowbreak esp32} e
\texttt{ls-praxedes/\allowbreak scanner-obd2/\allowbreak qsr5c99}), o serviço de
ingestão no \textit{backend} subscreve (\texttt{SUBSCRIBE}) utilizando o caractere
coringa de nível único (\texttt{ls-praxedes/\allowbreak scanner-obd2/+}), viabilizando
uma arquitetura \textit{multi-tenant} capaz de agregar múltiplos veículos
simultaneamente em um único trabalhador (\textit{worker}) de processamento.

Para mitigar na prática os vetores de escuta passiva (\textit{eavesdropping}) e de
intermediário (\acrshort{mitm}) caracterizados na Seção~\ref{sec:vulnerabilidades_mqtt},
a segurança do sistema foi projetada sob a arquitetura \textit{Zero Trust} (NIST SP
800-207), estabelecendo a distinção formal entre criptografia de canal em trânsito e
autorização de acesso \cite{rose2020,baars2018}. Enquanto o transporte sobre
\acrshort{tls} v1.2 com verificação matemática da cadeia X.509 assegura a
confidencialidade e a autenticidade da conexão, a arquitetura especifica para ambientes
de produção o bloqueio de conexões anônimas no \textit{broker} e a aplicação de Listas
de Controle de Acesso (\acrshort{acl}s, do inglês \textit{Access Control Lists}) regidas
pelo Princípio do Menor Privilégio, restringindo o perfil credenciado do veículo
exclusivamente à permissão de escrita (\texttt{write}) em seu identificador próprio e
reservando a permissão de leitura (\texttt{read}) em escopo coringa unicamente ao
serviço de ingestão do \textit{backend}.

Internamente, o serviço de aplicação foi estruturado sob o padrão de Arquitetura em
Camadas (\textit{Layered Architecture}) \cite{fowler2002}, combinado às diretrizes de
portabilidade e ciclo de vida da metodologia \textit{The Twelve-Factor App}
\cite{wiggins2011}:

\begin{itemize}
    \item \textbf{Fonte Única da Verdade (\acrshort{ssot}):} Em cumprimento ao Fator III
        (\textit{Config}) do \textit{The Twelve-Factor App}, centralizaram-se em módulo
        dedicado todas as constantes de infraestrutura (\texttt{MQTT\_BROKER},
        \texttt{MQTT\_PORT = 8883}, \texttt{MQTT\_TOPIC},
        \texttt{MQTT\_\allowbreak MOCK\_\allowbreak TOPIC}) e a resolução dinâmica de
        caminhos absolutos de sistema de arquivos via biblioteca \texttt{pathlib.Path}
        (\texttt{CA\_CERT\_PATH} e \texttt{DATABASE\_URL}), garantindo portabilidade
        entre sistemas operacionais e eliminando divergências de parametrização.
    \item \textbf{Ingestão Assíncrona e Orquestração \acrshort{asgi}:} O cliente de
        mensageria Eclipse Paho-MQTT é inicializado com contexto criptográfico
        \texttt{ssl.PROTOCOL\_\allowbreak TLSv1\_2} (\texttt{tls\_set}) e acoplado ao
        servidor \textit{web} FastAPI (executado sobre o servidor \acrshort{asgi}
        Uvicorn) por intermédio do gerenciador de contexto assíncrono
        \texttt{@asynccontextmanager lifespan(app)}. Na partida da \acrshort{api}, a
        invocação de \texttt{client.\allowbreak loop\_start()} instancia uma
        \textit{thread} de segundo plano não bloqueante para escuta contínua do
        \textit{broker} sem obstruir o laço de eventos principal de atendimento a
        requisições \acrshort{http}; no encerramento do processo, executam-se
        \texttt{client.\allowbreak loop\_stop()} e
        \texttt{client.\allowbreak disconnect()}, assegurando o desligamento gracioso
        (\textit{graceful shutdown}, Fator IX -- \textit{Disposability}).
\end{itemize}

Para proteger a integridade do banco de dados perante pacotes ruidosos ou inválidos, a
função de retorno (\textit{callback}) de recepção \texttt{on\_message} implementa um
mecanismo defensivo em três níveis: (i) captura de exceções de decodificação sintática
(\texttt{json.\allowbreak JSONDecodeError}) para descarte imediato de cargas binárias
truncadas; (ii) sanitização estrita de contrato e tipos de domínio via Pydantic v2 (por
meio das classes de esquema \texttt{TelemetryCreate} e
\texttt{Telemetry\allowbreak Response} com \texttt{from\_attributes = True}), atuando
como um \textit{firewall} de aplicação que rejeita pacotes malformados
(\texttt{pydantic.\allowbreak ValidationError}) antes da camada de persistência; e (iii)
bloco transacional \texttt{try...except...finally} sobre o mapeamento objeto-relacional
SQLAlchemy 2.0 \acrshort{orm}, executando \texttt{db.commit()} mediante sucesso ou
\texttt{db.rollback()} em caso de falha de entrada e saída, sucedido invariavelmente
pelo fechamento da sessão via \texttt{db.close()} para evitar vazamento de descritores
de arquivo.

Os registros validados são persistidos de maneira aderente às propriedades
\acrshort{acid} na tabela relacional \texttt{telemetry} do banco de dados SQLite, cujo
esquema compreende as colunas \texttt{id} (chave primária inteira indexada),
\texttt{rpm} (\texttt{Float}, não nulo), \texttt{temperature} (\texttt{Float}, não
nulo), \texttt{speed} (\texttt{Float}, anulável) e \texttt{timestamp} (\texttt{DateTime}
registrado em fuso horário UTC com precisão de microssegundos). Para consumo externo, a
camada de aplicação expõe os \textit{endpoints} RESTful \texttt{/api/v1/telemetry}
(\texttt{POST} e \texttt{GET} com suporte a paginação via parâmetros \texttt{skip} e
\texttt{limit}, ordenação cronológica decrescente
\texttt{Telemetry.\allowbreak timestamp.\allowbreak desc()} e Injeção de Dependência via
\texttt{Depends(get\_db)}). A corretude de toda a infraestrutura de servidor foi
verificada por uma suíte automatizada de testes unitários com o \textit{framework}
\texttt{pytest}, abrangendo testes de rotas via \texttt{TestClient} e testes isolados do
trabalhador de mensageria com objetos simulados
(\texttt{unittest.\allowbreak mock.\allowbreak MagicMock}), validando tanto a
persistência de pacotes íntegros quanto a rejeição determinística de cargas corrompidas
sem dependência de \textit{hardware} físico.

\section{Emulação veicular, visualização em tempo real e modelagem da camada de detecção
de anomalias}

Com o propósito de viabilizar testes de carga e validação contínua da esteira
telemétrica independentemente do veículo em rodagem, desenvolveu-se um emulador
sintético de dinâmica veicular em Python. O utilitário consome os parâmetros
centralizados de infraestrutura, estabelece sessão segura \acrshort{mqtts}
(\acrshort{tls} v1.2 na porta 8883) mediante validação da cadeia de certificação X.509 e
publica pacotes \acrshort{json} a cada $0{,}5\text{ s}$, simulando variações
estocásticas e regimes operacionais realistas de rotação, temperatura e velocidade.

Na camada de apresentação, a interface analítica interativa foi construída utilizando o
\textit{framework} Streamlit. A aplicação consome periodicamente o \textit{endpoint}
RESTful \texttt{GET /api/v1/telemetry}, estrutura as séries temporais recebidas em
objetos \textit{DataFrame} da biblioteca Pandas (convertendo os carimbos
\texttt{timestamp} para objetos temporais e assegurando a indexação cronológica
crescente para plotagem fiel da evolução temporal da esquerda para a direita), computa
variações instantâneas entre leituras consecutivas ($\Delta\text{RPM}$,
$\Delta T_{\text{ect}}$ e $\Delta v$) e a grandeza derivada de aceleração cinemática
instantânea ($a = \Delta v / \Delta t$), além de monitorar o estado de conectividade
ativa do veículo. Para eliminar o efeito visual de recarregamento integral da página
(\textit{flicker}) e a latência de redesenho, implementou-se a atualização seletiva da
árvore de interface por meio de um contêiner dinâmico (\texttt{st.empty()}) realimentado
em laço de $0{,}5\text{ s}$ (\texttt{auto\_refresh}), renderizando cartões de
indicadores-chave (\textit{metrics}) e três gráficos de linha independentes para
rotação, temperatura do arrefecimento e velocidade, o que previne a distorção visual de
escala decorrente da sobreposição de grandezas de ordens de magnitude distintas.

No escopo de desenvolvimento incremental do projeto, enquanto os subsistemas de
aquisição embarcada no barramento \acrshort{can}, transporte criptografado
\acrshort{mqtts}, ingestão validada, persistência relacional \acrshort{acid} e
visualização em tempo real encontram-se integralmente implementados e operacionais, a
camada de inteligência artificial é apresentada neste trabalho em nível de especificação
arquitetural e modelagem analítico-matemática projetada para integrar a plataforma como
etapa subsequente do fluxo de dados. Retomando o ciclo metodológico \acrshort{crispdm}
fundamentado na Seção~\ref{sec:crispdm_iforest} \cite{wirth2000,geron2022}, a modelagem
do \textit{pipeline} para detecção não supervisionada de anomalias operacionais
estrutura-se em três estágios aplicados às grandezas do motor:

\begin{enumerate}
    \item \textbf{Compreensão do Negócio e dos Dados (\textit{Business and Data
        Understanding}):} Caracterização do acoplamento termodinâmico e cinemático
        nominal do motor a combustão interna, compreendendo a faixa de estabilização
        térmica do líquido de arrefecimento (tipicamente entre
        $85\text{ }^{\circ}\text{C}$ e $105\text{ }^{\circ}\text{C}$ após o aquecimento
        inicial) e a relação de proporcionalidade entre rotação do motor e velocidade
        veicular em cada marcha de transmissão.
    \item \textbf{Preparação dos Dados (\textit{Data Preparation}):} Extração das séries
        históricas persistidas no banco relacional, remoção de registros incompletos,
        filtragem de transientes espúrios de ignição e estruturação da matriz de
        atributos multivariada $\mathbf{X}$, na qual cada observação temporal $i$ é
        representada pelo vetor de características
        $\mathbf{x}_i = [\text{RPM}_i, T_{\text{ect},i}, v_i]^T$.
    \item \textbf{Modelagem Analítica (\textit{Isolation Forest}):} Formulação do modelo
        não supervisionado baseado no princípio de particionamento recursivo estocástico
        em árvores binárias de isolamento (\textit{iTrees}) exposto na
        Seção~\ref{sec:crispdm_iforest}.
\end{enumerate}

A quantificação formal do grau de anormalidade de cada vetor telemétrico $\mathbf{x}$ em
uma amostra de $n$ instâncias é determinada pelo escore de anomalia $s(\mathbf{x}, n)$,
definido por \citeonline{liu2008} conforme a Equação~\eqref{eq:iforest}:

\begin{equation}\label{eq:iforest}
s(\mathbf{x}, n) = 2^{-\frac{E(h(\mathbf{x}))}{c(n)}}
\end{equation}

\noindent onde $h(\mathbf{x})$ representa o comprimento de caminho
(\textit{path length}, número de arestas percorridas desde o nó raiz até o nó folha de
isolamento) da amostra $\mathbf{x}$, $E(h(\mathbf{x}))$ denota a média aritmética de
$h(\mathbf{x})$ através de todas as árvores de isolamento da floresta e $c(n)$
corresponde ao fator de normalização dado pelo comprimento médio de uma busca
malsucedida em uma Árvore Binária de Busca (\textit{Binary Search Tree}) com $n$ nós.
Conforme expresso na Equação~\eqref{eq:iforest}, pontos operacionais situados em regiões
esparsas da distribuição multivariada, tais como episódios de sobreaquecimento térmico
severo ou elevação abrupta de rotação desacoplada da carga cinemática, requerem poucas
partições para serem isolados, resultando em $E(h(\mathbf{x}))$ reduzido e conduzindo o
escore ao limite assintótico $s(\mathbf{x}, n) \to 1$. Tais instâncias são classificadas
pelo modelo com o rótulo \texttt{-1} (anomalia operacional), em contraste às amostras de
regime estacionário normal rotuladas como \texttt{+1} ($s(\mathbf{x}, n) \le 0{,}5$).

\section{Apoio instrumental de IAG na verificação de lacunas e revisão}

Em conformidade com as diretrizes institucionais estabelecidas pela Resolução Normativa
nº 17/2025 e pela Instrução Normativa nº 06/2026 da Pró-Reitoria de Pós-Graduação e
Pesquisa da Universidade Federal de São Paulo \cite{unifesp_iag2026}, registra-se que,
ao longo do desenvolvimento do \acrshort{tcc} II durante o segundo semestre de 2026,
utilizou-se de forma complementar a ferramenta de Inteligência Artificial Generativa
(\acrshort{iag}) Google Gemini (integrada ao ambiente de desenvolvimento) estritamente
como instrumento de apoio à revisão e à verificação exploratória de lacunas
(\textit{gap analysis}). O emprego da ferramenta circunscreveu-se a três atividades
auxiliares: (i) verificação de lacunas técnicas e casos de borda
(\textit{gap checking}), auxiliando o autor na inspeção de potenciais pontos cegos de
concorrência entre tarefas no \textit{firmware} embarcado e no levantamento de cenários
de exceção para a cobertura de testes unitários do \textit{backend}; (ii) consulta
remissiva a normas técnicas, agilizando a localização de cláusulas específicas em
padrões de engenharia (tais como as regras de subscrição com caracteres coringa do
protocolo OASIS \acrshort{mqtt}, diretrizes de encapsulamento MISRA C:2012 e normas de
formatação bibliográfica ABNT); e (iii) revisão linguística, gramatical e de organização
estrutural das seções textuais previamente concebidas pelo autor.

Ressalta-se com rigor que a concepção arquitetural, a formulação lógica, a implementação
integral do código-fonte, a condução experimental e a interpretação dos resultados
constituíram atividades intelectuais exclusivas e intransferíveis do autor. Para
assegurar a fidedignidade científica do trabalho, todo apontamento de lacuna ou sugestão
textual oriundo da ferramenta foi submetido a um protocolo sistemático de validação
humana e empírica em três etapas antes de sua consolidação: (1) confronto direto com a
literatura primária, verificando a exatidão física e normativa perante os
\textit{datasheets} oficiais dos fabricantes (Espressif Systems e Texas Instruments) e
as especificações internacionais (ISO 15765-4, \acrshort{sae} J1979 e IETF RFC 5246);
(2) verificação determinística de \textit{software}, submetendo o código embarcado à
compilação estrita sem avisos no compilador \texttt{xtensa-esp32-elf-gcc} e a camada de
servidor à bateria automatizada de testes unitários com \texttt{pytest}; e (3) validação
empírica em bancada e no veículo alvo, comprovando o determinismo temporal e funcional
da solução no microcontrolador ESP32.

\chapter{Planejamento e Cronograma de Execução}

O cronograma foi estruturado para concentrar as atividades de pesquisa, planejamento e
projeto no \acrshort{tcc} I, estabelecendo uma base teórica e metodológica sólida
(Tabela~\ref{tab:cronograma_tcc1}). As atividades práticas de desenvolvimento, montagem,
testes e análise de resultados estão alocadas para o \acrshort{tcc} II
(Tabela~\ref{tab:cronograma_tcc2}), focando na implementação e validação da solução
proposta.

\begin{table}[!htb]
\centering
\caption{Cronograma de Atividades -- \acrshort{tcc} I (2025/1)}
\label{tab:cronograma_tcc1}
\begin{tabular}{|l|c|c|c|c|c|c|}
\hline
\multicolumn{1}{|c|}{\multirow{2}{*}{\textbf{Foco Teórico e Planejamento}}}
    & \multicolumn{6}{c|}{\textbf{Meses}} \\ \cline{2-7}
\multicolumn{1}{|c|}{} & \textbf{Fev} & \textbf{Mar} & \textbf{Abr} & \textbf{Mai}
    & \textbf{Jun} & \textbf{Jul} \\ \hline
1. Revisão Bibliográfica Aprofundada & \cellcolor[HTML]{C0C0C0}
    & \cellcolor[HTML]{C0C0C0} & \cellcolor[HTML]{C0C0C0} & & & \\ \hline
2. Definição da Arquitetura do Sistema & & & \cellcolor[HTML]{C0C0C0}
    & \cellcolor[HTML]{C0C0C0} & & \\ \hline
3. Detalhamento do Projeto & & & \cellcolor[HTML]{C0C0C0} & \cellcolor[HTML]{C0C0C0}
    & \cellcolor[HTML]{C0C0C0} & \\ \hline
4. Estudo de Algoritmos para Análise & & & & & \cellcolor[HTML]{C0C0C0}
    & \cellcolor[HTML]{C0C0C0}\\ \hline
5. Compra e Entrega de Componentes & & & & &  & \cellcolor[HTML]{C0C0C0} \\ \hline
6. Redação da Monografia (Parte I) & \cellcolor[HTML]{C0C0C0} & \cellcolor[HTML]{C0C0C0}
    & \cellcolor[HTML]{C0C0C0} & \cellcolor[HTML]{C0C0C0} & \cellcolor[HTML]{C0C0C0}
    & \cellcolor[HTML]{C0C0C0} \\ \hline
\end{tabular}
\legend{Fonte: Elaborado pelo autor (2026).}
\end{table}

\begin{table}[!htb]
\centering
\caption{Cronograma de Atividades -- \acrshort{tcc} II (2026/2)}
\label{tab:cronograma_tcc2}
\begin{tabular}{|l|c|c|c|c|c|}
\hline
\multicolumn{1}{|c|}{\multirow{2}{*}{\textbf{Foco Prático e Desenvolvimento}}} &
\multicolumn{5}{c|}{\textbf{Meses}} \\ \cline{2-6}
\multicolumn{1}{|c|}{} & \textbf{Ago} & \textbf{Set} & \textbf{Out} & \textbf{Nov}
    & \textbf{Dez} \\ \hline
7. Montagem e Teste do Hardware & \cellcolor[HTML]{C0C0C0} & & & & \\ \hline
8. Dev. Leitor de Protocolo & \cellcolor[HTML]{C0C0C0} & \cellcolor[HTML]{C0C0C0}
    & \cellcolor[HTML]{C0C0C0} & & \\ \hline
9. Dev. Armazenamento e Visualização & & \cellcolor[HTML]{C0C0C0}
    & \cellcolor[HTML]{C0C0C0} & & \\ \hline
10. Coleta e Validação inicial de dados & & \cellcolor[HTML]{C0C0C0}
    & \cellcolor[HTML]{C0C0C0} & & \\ \hline
11. Dev. Módulo de Análise Inteligente & & & \cellcolor[HTML]{C0C0C0}
    & \cellcolor[HTML]{C0C0C0} & \\ \hline
12. Integração e Testes Finais & & & & \cellcolor[HTML]{C0C0C0}
    & \cellcolor[HTML]{C0C0C0} \\ \hline
13. Redação da Monografia (Parte II) & \cellcolor[HTML]{C0C0C0}
    & \cellcolor[HTML]{C0C0C0} & \cellcolor[HTML]{C0C0C0} & \cellcolor[HTML]{C0C0C0}
    & \cellcolor[HTML]{C0C0C0} \\ \hline
14. Preparação para a Defesa Final & & & & &  \cellcolor[HTML]{C0C0C0} \\ \hline
\end{tabular}
\legend{Fonte: Elaborado pelo autor (2026).}
\end{table}

% =========================================================================
% ELEMENTOS PÓS-TEXTUAIS (abnTeX2 + abntex2cite)
% =========================================================================
\postextual
\bibliography{bibliography}

\end{document}